% =====================================================================
%  GENERATED; DO NOT EDIT BY HAND.  Rebuild: sh ebook-method/etiske-undersoegelser/rebuild.sh
%  Harald Høffding, Etiske Undersøgelser.  Kjøbenhavn: P. G. Philipsens Forlag, 1891.  [VI] + 83 pp.; three treatises:
%  Om Muligheden af en filosofisk Etik (pp. 1-23), Om Velfærdsprincipet (pp. 24-41), Forholdsloven i Etiken (pp. 42-83).
%  Diplomatic transcription (HOEFFDING-QUEUE #6), begun 2026-10-16.
%  ACCURACY (2026-10-16).  Scan: Google's scan of the Wisconsin copy (Internet Archive etiskeundersgel00hfgoog), 600 ppi
%  bitonal, [VI] + 83 pp., complete by folio; no page damaged by the measure of mindre-arbejder-2 (low-confidence line-edge
%  words: max 6, threshold 8), gutter not cut (33-40 pt of white).  The NB copy (8df6665866cef49a09cc7815de97c99b) could not be
%  fetched (NB's server returned 404/500 for its images and ALTO, 2026-10-16): ONE copy, two readings.
%  Readings: our tesseract (tessdata_best dan, 300 dpi; own non-words 0.33 % against ABBYY's 1.15 %) as the draft, the
%  Archive's ABBYY as the witness; collated word by word (ebook-method/collate/check.py): 588 disagreements, 54 settled by
%  check.py, 189 by rules.py, the rest at the image.  Blind compact crop sheets (READER-BRIEF.md: 549 word spots in 535 crops
%  + 33 paragraph starts, 28 planted controls, 28 right) covered every open check, an audit sample of 12 per rule class
%  (one class, markplaced, concerns note marks only and is superseded by the note reading; layerfault wrong once, forudsé),
%  the unattested-word sweep, specks and an ø/umlaut sweep; then whole-line crops (z1) for what the word crops could not
%  settle (every line-end dash tesseract alone read proved to be nothing).  Decisions in decisions.tsv (rule:/hand: notes),
%  word fixes in fixes1.tsv.  The 53 footnotes were read whole by two note readers and set from their reading at the marks
%  they read (noteplace.py); their text was diffed against the collated draft's (notesdiff.py): 26 differences, all the
%  draft's except one reader's normalisation (Vierterljahrsschrift, p. 20, settled at 600 ppi).  Structure, front matter, the
%  fraction (p. 55) and the Greek (p. 57) from a structure reader, the Greek also read in the blind crops (two readings agree).
%  Letterspacing MEASURED (spacemeasure.py: white runs bimodal, GAP 5; 153 runs), checked on a sample of 30 runs + 12 lines
%  (no false run; one too wide, one short word missed) and set by spaceapply.py.
%  Blind page check, STRATIFIED: 3 pages of the 15 less clean (>= 3 low-confidence edge words: pp. 25, 57, 58) and 3 clean
%  (pp. 13, 26, 64), seed 2026101606, 1,686 words: 0 word errors from the collation; 1 fault introduced by a patch (a %
%  comment after the Greek hid 7 words on p. 57) and 2 letterspacing faults (p. 57 Aristoteles not set; p. 13 ham set):
%  all mended.  Rates: word errors 0.59 / 1000 (95 % upper 2.81), emphasis 1.19 / 1000 (upper 3.73).
%  RESIDUE, stated honestly: (1) letterspacing is measured, not read: short words beside a run and runs at the end of a
%  line may be missed or over-set (~1 per 1000 words by the check); 6 measured runs could not be placed (p. 6 forkastes,
%  in notes Körner, Darwin, Jhering are set from the note readers' own marking); (2) the two OCRs read the same images, so
%  a fault both share where the corpus attests the word is invisible to the collation (the translators' odd spots are
%  the next reading); (3) the Greek accents of p. 57 rest on two readings at ~300 dpi (ἓν as both read; Aristotle ἕν).
%  THE TRANSLATORS' ODD SPOTS (52, from four translators) were all looked at in 600 ppi line crops of the scan (main
%  session): 12 faults of the transcription mended (specks read as , . ; and a stop; a space before «; praktisk-populære;
%  juridisk-teologiske; tesseract junk 5% and the folio 838; the p. 6 note repeated in the body; p. 82 a run wrongly set
%  letterspaced; p. 33 Tillempelsen, a comma after all) -- fixes1.tsv („odd“ rows), patches.NOTECUT; the rest are the print's.
%  Misprints kept as printed, marked PRINTED AS IS at the line (patches.MISPRINTS, 17): the print's weak f reads as t in
%  „af“ six times (pp. 8, 9, 16, 40, 43, 54: „at“ for „af“), doubled words at line turns (pp. 16, 32, 65, 69, 82), a
%  missing word (pp. 39, 73), vidde (48), Crimonologie (74), sukjektivt (57), a comma ending a paragraph (38).  In the
%  notes, as printed and not marked there: p. 20 Vierterljahrsschrift; p. 43 „Lectures ou Jurisprudence“; p. 62
%  Gefängsnisswesens and Tardé (p. 77 Tarde); p. 70 „affectus coërceri nec tolli“ (Spinoza: nec coërceri).  Printed and
%  right: hints (p. 50, genitive of hin), Ildlænderne (pp. 32-33, the Fuegians), Shakesperes (p. 11), Arbeide (p. 79).
%  The print has no å: tesseract's å for ä mended (Vorgänge, Clärchen, Aktenmässige).
%  Conventions: printer's errors transcribed as printed and logged in a % comment at the site; \opage{N} = printed page N
%  begins; a word broken over a page turn is split with %; footnotes as printed, marked *), **), ***) and numbered per
%  page (\footnote[k]).  Italics \textit; letterspacing \emph.  The ø rule (CLAUDE.md) applies.
% =====================================================================
\documentclass[12pt,a4paper,openany]{book}
\usepackage[utf8]{inputenc}
\usepackage[T1]{fontenc}
\usepackage{libertinus}
\usepackage{libertinust1math}
\usepackage[danish]{babel}
\usepackage{textalpha}   % the Greek on p. 57
\usepackage[protrusion=true,expansion=true]{microtype}
\usepackage{setspace}
\usepackage[margin=1.2in]{geometry}
\usepackage{marginnote}
\usepackage{hyperref}
\hypersetup{pdftitle={Etiske Undersøgelser}, pdfauthor={Harald Høffding}, hidelinks}
\setstretch{1.3}
\renewcommand{\thefootnote}{%
  \ifcase\value{footnote}\or *)\or **)\or ***)\or \dag)\or \dag\dag)\or \dag\dag\dag)\else \arabic{footnote})\fi}   % p. 28: six notes
\makeatletter\@addtoreset{footnote}{page}\makeatother
\newcommand{\opage}[1]{\marginnote{\footnotesize\textit{[#1]}}}
\newcommand{\essay}[2][]{\chapter*{#2}\addcontentsline{toc}{chapter}{#1}}   % [toc title]{title as printed}
\pdfstringdefDisableCommands{\renewcommand{\footnote}[2][]{}\renewcommand{\opage}[1]{}}
\begin{document}

\frontmatter
% Front matter (PDF 7 title, PDF 9 Forord, PDF 11 Indhold), set from the structure reading (struct-S1) checked against the
% layer.  The scan's library stamps and pencil (A 1289, 623724) are not print.
\begin{titlepage}\centering
{\large\textsc{Harald Høffding.}}\\[1ex]\rule{3em}{0.4pt}\\[3ex]
{\Huge ETISKE UNDERSØGELSER.}\\[4ex]
\fbox{\begin{tabular}{c}\small\textbf{OM MULIGHEDEN AF EN FILOSOFISK ETIK.}\\\small\textbf{OM VELFÆRDSPRINCIPET.}\\\small\textbf{FORHOLDSLOVEN I ETIKEN.}\end{tabular}}\par\vfill
{\small KØBENHAVN.}\\[0.5ex]{\small\textbf{P. G. PHILIPSENS FORLAG.}}\\[0.5ex]{\footnotesize\textbf{THIELES BOGTRYKKERI.}}\\[1ex]{\small 1891.}
\end{titlepage}
\chapter*{FORORD.}
Af de tre Afhandlinger, der udgøre dette Skrift, lægger jeg selv størst Vægt paa den tredie, i hvilken jeg har søgt at
gennemføre en i min Etik (noget mere i den tyske end i den danske Udgave) antydet Idé, der turde være af nogen Betydning
for den etiske Tænkning. Denne Afhandling har allerede været trykt i første Hefte af »International Journal of Ethics«
(Efteraaret 1890). --- I de to første Afhandlinger har jeg taget nogle Spørgsmaal, som angaa Etikens første Principer, op
til fornyet Behandling. De kunne betragtes enten som Indledning til eller som Supplement til min Etik (især Kap.~I, III og
VII). Den anden Afhandling (Om Velfærdsprincipet) trykkes omtrent samtidigt hermed i det amerikanske Tidsskrift »The
Monist« som Del af en længere Artikel. ---

\noindent\hspace*{2em}\textit{Januar 1891.}\hfill\textbf{Harald Høffding.}
\chapter*{INDHOLD.}
\begin{flushright}\small\textbf{Side}\end{flushright}
\noindent\begin{tabular}{@{}p{0.86\textwidth}@{\hspace{0.5em}}r@{}}
Om Muligheden af en filosofisk Etik\dotfill & 1\\
Om Velfærdsprincipet\dotfill & \textbf{24}\\
Forholdsloven i Etiken\dotfill & \textbf{42}\\
\end{tabular}
\mainmatter
% --- p. 1 ---
\setcounter{footnote}{0}%
\essay[Om Muligheden af en filosofisk Etik]{Om Muligheden af en filosofisk Etik.}
\opage{1}
1. At prædike Moral er let, at begrunde Moral er
vanskeligt, har en Filosof sagt. Denne Sætning er rigtig,
dersom man ved at prædike Moral blot forstaar at udtale
de i det givne Samfund gængse Forestillinger om Godt
og Ondt. Men hvis man ved at prædike Moral forstaar
at henvende sig paa en saadan Maade til de enkelte
Individer, at disses Karakter derved faar Motiver og
Tilskyndelser til at udvikle sig ikke blot i de gængse
etiske Forestillingers Aand, men i sande og gyldige Idealers
Aand, saa er det ikke blot den vanskeligste og den
højeste Kunst af alle, men den vil ogsaa nødvendigvis
føre over til hvad hin Filosof erkendte som en stor Vanskelighed,
nemlig at give en Begrundelse af etiske Domme.

Etiske Domme, Domme om Godt og Ondt, opstaa
fra først af som uvilkaarlige Følelsesudbrud. De udtale
den Værdi, der tillægges menneskelige Handlinger eller
Institutioner. I ethvert Samfund er der et vist Indbegreb
af etiske Domme, hvis Gyldighed anerkendes. De ligge
i Luften, idet det enkelte Individ tilegner sig og anvender
disse Domme som umiddelbart givne og berettigede, lige
som det arver Racens Instinkter og Traditioner i det
Hele. Dette Indbegreb af etiske Domme udtrykker hvad
man har kaldet \emph{den positive Moralitet}. Denne har,
% --- p. 2 ---
\setcounter{footnote}{0}%
\opage{2}lige som positiv Religion og positiv Lovgivning, en historisk
Karakter. Den har udviklet sig og udvikler sig
stadigt under Vexelvirkning med Racens eller Samfundets
Skæbne og Livsforhold. Dens Indhold er en Kreds af
Forbilleder, Leveregler og Bud, der i lange Tider gaa
fra Slægt til Slægt som »uskrevne Love« og ofte staa
for den Enkelte med hemmelighedsfuld Myndighed. Det
enkelte Individ bøjer sig saa meget mere for dem, som
de forkyndes og hævdes af Autoriteter, Personligheder,
hvis fysiske eller aandelige Magt vækker dets Frygt, dets
Ærefrygt, dets Tillid eller dets Kærlighed, og over for
hvilke ingen Trang til Kritik føles, i hvert Tilfælde ingen
Kritik let kommer til Orde. Saadanne Autoriteter ere
Forældre, Lærere, Præster, Fyrster, Lovgivere, Religionsstiftere.
--- Den positive Moralitet kunde tænkes fremstillet
i sin Sammenhæng. Vi vilde da faa en Fremstilling
af de hos et vist Folk og paa en vis Tid gældende etiske
Domme. Den teologiske Etik og den positive Retsvidenskab
vilde være Dele af en saadan positiv\footnote[1]{Ordet „positiv“ betegner i hele denne Sammenhæng (positiv Moralitet, positiv Religion, positiv Lovgivning, positiv Etik) ikke Modsætning til „negativ“. Heller ikke bruges det i den Betydning, i hvilken Auguste Comte tager det, naar han kalder sin Filosofi „la philosophie positive“ eller Positivismen. Hos Comte er „positif“ det Samme som faktisk, virkeligt givet i Naturen. Derimod betegner Ordet „positiv“ her, hvad der er sat i bestemt \emph{historisk} Form, som udtrykkeligt fastslaaet Bestemmelse eller \emph{Tradition}.} Etik (saaledes
kunde man maaske kalde den videnskabelige Fremstilling
af den positive Moralitet), idet den teologiske Etik
blot angaar de etiske Domme, der kunne udledes af en
vis religiøs Tro, og den positive Retsvidenskab kun omfatter
de Regler og Bud, der udtrykkeligt ere formulerede
i Landets juridiske Love. Den positive Etik vilde tillige
% --- p. 3 ---
\setcounter{footnote}{0}%
\opage{3}indeholde de Regler og Bud, som den offentlige Mening,
Sæd, Skik, Tradition og Undervisning indskærpe.

Men naar Frygten for eller Tilliden til de historisk
givne Autoriteter svinder, eller naar af andre Aarsager
Trangen til en dybere gaaende etisk Erkendelse vaagner,
kan man ikke blive staaende ved den positive Moralitet.
Der reiser sig et Spørgsmaal om selve Autoriteternes
Berettigelse, om de første Principer, der betinge Gyldigheden
af alle etiske Domme. Have etiske Domme ikke
Gyldighed, uden for saa vidt de hævdes af Autoriteter?
Det er det store Spørgsmaal, ved hvis Fremsættelse \emph{den
filosofiske Etik} bliver til. Dennes Opgave maa være
at vise, at man uden at forudsætte Autoriteternes absolute
Gyldighed, kan finde et Grundlag for etiske Domme
i den menneskelige Natur. Lige som den filosofiske Erkendelsesteori
angaar de første Principer for den teoretiske
Verdensopfattelse, som grunder sig paa den
menneskelige Natur, saaledes angaar den filosofiske Etik
de første Principer for den \emph{Vurdering} af menneskelige
Handlinger, som grunder sig paa selve den menneskelige
Natur.

2. Forholdet mellem den filosofiske Etik og den
positive Moralitet behøver ikke at være fjendtligt. Den
menneskelige Natur, i hvilken Filosofien søger et Grundlag
for etiske Domme, aabenbarer sig ogsaa i Historien
og faar bestemmende Indflydelse paa Udviklingen af den
positive Moralitet. De Idealer og Bud, som udgøre dennes
Indhold, ere ikke opstaaede uden i Kraft af den menneskelige
Natur, saaledes som denne under de givne Forhold
virkede. Selv i den strengeste Autoritetstro lægger
derfor Menneskenaturen sig for Dagen. Ingen Autoritet
kan være rent ydre: Autoriteten er kun Autoritet ved
den Følelse, den formaar at vække. Troen paa Autoriteten
hænger derfor stedse sammen med Menneskenes  hele
% --- p. 4 ---
\setcounter{footnote}{0}%
\opage{4}Følelsesliv og gennem dette igen med de Erfaringer, som
Menneskene gøre under deres Kamp for Tilværelsen.
Ved saaledes at gaa tilbage til den positive Moralitets
Sammenhæng med den menneskelige Natur, ser man dens
Slægtskab med den filosofiske Etik. For en Del har den 
filosofiske Etik kun den Betydning \emph{at bringe vore
etiske Dommes Sammenhæng med vor Natur og
vore Livsvilkaar til fuld Bevidsthed}, som jo i
det Hele Filosofien kan siges at være den menneskelige
Aands Selverkendelse i teoretisk og praktisk Henseende.


Men foruden dette Forhold mellem filosofisk Etik og
positiv Moralitet, der fremkommer ved, at denne ikke
har den samme Trang som hin til at gaa tilbage til det
første Grundlag, kan der endnu finde tre forskellige Forhold
Sted mellem dem.

Den filosofiske Etik kan ud fra den klarere Forestilling
om de første Principer føres til at fordre, at der
skal \emph{arbejdes videre} væsentligt i samme Aand som
inden for den positive Moralitet, men ud fra mere omfattende
Synspunkter. Saaledes arbejder den nyere Tids
Humanitetsfølelse videre paa Grundlag af det Menneskekærlighedens
Princip, som især den kristelige Etik hævdede,
som den græske Etik gjorde gældende inden for
de ved Nationalfølelsen satte Grænser, og som allerede
virkede i den gammelindiske Familieetik i den Følelse,
der omfattede Alle, der samledes om samme Arne. Det
gælder her blot Principets \emph{Udvidelse} ud over Konfessionens,
Nationens og Familiens Grænser\footnote[1]{Smlgn. min „Etik“. København 1887. p. 326---331.}.

Medens det her gælder Udvidelse og Gennemførelse,
kan det i andre Tilfælde være nødvendigt at indføre en
helt anden Motivering og Begrundelse, for at Anerken%
% --- p. 5 ---
\setcounter{footnote}{0}%
\opage{5}delse og fortsat Anvendelse af det historisk Givne kan
blive mulig. Det historisk Givne i Handlemaade eller i
Samfundsforholdenes Ordning anerkendes og benyttes da
ikke af det Motiv, der oprindeligt har fremkaldt det, men
af et helt andet Motiv. Der finder med andre Ord en
Motivforskydning Sted. --- Af Arnens Hellighed udviklede
sig hos de gamle Indere Asylretten, der særligt
blev af Betydning ved Arnestederne i Templet og i
Kongens Gaard\footnote[1]{Smlgn. \emph{Leist}: Alt-arischen Jus gentium. Leipzig 1889. p. 20 f. 87. --- F. \emph{Pollock}: The King’s Peace (Oxford Lectures. London 1890). p. 70. 77.}; denne Asylret blev igen en Form,
under hvilken Fred kunde stiftes og Benaadning øves;
og heraf udsprang dels Statsmagtens Haandhævelse af
Ret og Orden, dels ogsaa den humanere Behandling af
Forbrydere og i det Hele af dem, der ere i Konflikt med
det bestaaende Samfund. --- Et andet Exempel er den
private Ejendom, der historisk har udviklet sig under
Indflydelse af Erhvervsdrift, Virkelyst og Magtfølelse, men
som Samfundet anerkendte og regulerede i det Enkelte
for Fredens og Ordenens Skyld, og som nu i etiske og
nationaløkonomiske Systemer kun anerkendes, for saa
vidt den fremmer den Enkeltes Selvvirksomhed i Samfundets
Tjeneste\footnote[2]{Smlgn. min „Etik“. p. 264---266.}. --- Motivforskydningen har en oplysende
Parallel i den Maade, hvorpaa et og samme
Organ hos forskellige Arter af levende Væsener kan udøve
forskellige Funktioner, fordi det er afpasset efter forskellige
Forhold. Forlemmerne ere saaledes Griberedskaber
og Føleredskaber hos Mennesket, Gangredskaber hos
Hesten, Flyveredskaber hos Fuglen, Svømmeredskab hos
Hvalen. En og samme Grundform fremtræder omdannet
paa forskellig Maade under de forskellige Livsforhold.
% --- p. 6 ---
\setcounter{footnote}{0}%
\opage{6}Omvendt kunne forskellige Organer komme til at udføre
samme Funktion, naar levende Væsener af forskellig Art
for en Del have samme Livsbetingelser. Det er saaledes
højst forskellige Dele af Organismen, der hos Pattedyr,
Fiske, Krebsdyr og Insekter fungere som Aandedrætsorganer;
Fuglenes og Insekternes Vinger skyldes Udvikling
af helt forskellige Legemspartier\footnote[1]{Smlgn. \emph{Gervais}: Élements de zoologie. Paris 1871. p. 49 f. 54 ff. \emph{Darwin}: Om Arternes Oprindelse. p. 507. 519. \emph{Körner}: Die logischen Grundlagen der Klassifikation der Organismen. (Wundt’s Philosophische Studien. II.). p. 242.}. Motivforskydningen
paa det psykologisk-etiske Omraade svarer til
denne Omdannelse og Tillempelse paa det biologiske Omraade.
Ogsaa paa det biologiske Omraade er det jo i en
vis Forstand Motivet, der ændres, idet de samme Organer
komme til at tjene forskellige Formaal eller forskellige
Organer til at tjene samme Formaal. I Etiken kunne
Handlinger og Institutioner saaledes faa en anden Betydning
og derfor en anderledes motiveret Anerkendelse
end forhen. Lige saa lidt som man fra den »adaptive«
(ved Tillempelse opstaaede) Lighed mellem Sommerfugle
og egentlige Fugle tør slutte til et virkeligt Slægtskab
mellem dem, lige saa lidt kan man fra en Handlings
eller en Institutions nuværende etiske Betydning slutte
til dens historiske Oprindelse og Udvikling. Og lige saa
omvendt: fra den historiske Opstaaen kan der ikke uden
videre sluttes til Værdien. --- Motivforskydningens Fænomen
egner sig derfor godt til at belyse Forholdet mellem
den historiske og den filosofiske Behandling af Etiken\footnote[2]{Se herom min „Etik“ p. 8---10.}.

Endelig kan Forholdet dog ogsaa være det, at den
positive Moralitets Grundsætninger og Bud helt maa forkastes,
og nye Grundsætninger stilles op i Stedet. Da
er Indholdet altsaa negativt, ikke begrundende, fort%
% --- p. 7 ---
\setcounter{footnote}{0}%
\opage{7}sættende eller omdannende. Et saadant Forhold indtager
den filosofiske Etik til Askesen (for saa vidt denne erklærer
Lystfølelse i og for sig for ond og vil være mere
end Øvelse i fornøden Selvbeherskelse), til Intolerancen,
til Slaveriet og til andre barbariske Institutioner. ---

3. Vi vende nu tilbage til det Spørgsmaal, om der
i den menneskelige Natur gives et fast Grundlag for
etisk Vurdering, uafhængigt af de historisk givne Autoriteters
Indflydelse. Af dette Spørgsmaals Besvarelse afhænger
Muligheden af en filosofisk Etik.

Al Videnskab har tilsidst sit Grundlag i den menneskelige
Natur. Saaledes forudsætter Logiken Tænkeevnen
(Sammenligningsevnen, Evnen til at finde Lighed
og Forskel); Geometrien forudsætter desuden Evnen til
Rumsanskuelse; Fysiken og Kemien forudsætte desuden
endnu de forskellige ydre Sanser (Berøringssans, Syn,
Hørelse, Smag, Lugt), o. s. v. Hvad forudsætter nu Etiken
for sit Vedkommende?

Ved etiske Domme virker der et Element med, der
ikke spiller nogen Rolle ved rent logiske eller naturvidenskabelige
Domme. Det drejer sig ved dem om Mere end
at fastslaa Identitet og Forskel eller at paavise en Sammenhæng
mellem Fænomener. Beskaffenheden af en Handling
eller en Institution kan kun fremkalde en etisk Dom,
en Dom om Godt og Ondt, naar vedkommende Handling
eller Institution vækker en Følelse af Lyst eller Ulyst
hos os. Dersom vi vare blot anskuende, sansende og
sammenlignende Væsener, uden Evne til at føle Lyst og
Ulyst, vilde saadanne Domme ikke opstaa. Enhver Dom
opstaar ved Analyse af en given Totalopfattelse\footnote[1]{Se min „Formel Logik“. 2 Udg. p. 13.}. Det
Ejendommelige ved de etiske Domme er, at de analysere,
i bevidst Form udtale den Følelse, Synet af eller Fore%
% --- p. 8 ---
\setcounter{footnote}{0}%
\opage{8}stillingen om en Handling vækker hos os. Der forudsættes
Evne til at opfatte Handlingen og gøre sig dens
Karakter tydelig (ved Sammenligning); men dertil kommer
endnu Følelseselementet, den vækker, som en tredie Faktor
ved Siden af Kendskabet til de faktiske Forhold og
Evnen til at tænke over Handlingen. Erfaring godtgør
jo ogsaa, at etiske Diskussioner kunne bringes til at
maatte stanse paa tre Maader, nemlig ikke blot naar en
af Parterne overbevises om en Selvmodsigelse eller om
en faktisk Fejltagelse,\footnote[1]{Se „Formel Logik“. p. 24 f.} men ogsaa naar det er helt forskellige
Arter af Følelse, der hos de stridende Parter
bestemme de etiske Domme. De etiske Uenigheder, der
kun angaa Konsekventser eller faktiske Omstændigheder,
ere ikke af afgørende Natur. Anderledes forholder det
sig, hvor det er en forskellig Følelse, der bestemmer
Vurderingen, --- hvor \emph{Vurderingsmotivet} altsaa er
forskelligt.

Dersom nu Følelserne varierede i det Uendelige, ---
dersom der ikke var visse typiske Former for Følelse,
som med en vis Fasthed gøre sig gældende hos de enkelte
Individer hvert for sig og paa samme Maade hos flere
forskellige Individer, saa vilde ingen filosofisk Etik være
mulig. Men saaledes er det ikke. Vi finde hos hvert
normalt Individ en Følelse, som er den herskende hos
ham, i det Mindste i en vis Periode, saaledes at den udgør
Tyngdepunktet i hans Natur og vil bestemme hans
vurderende Domme. Det er den vigtigste Bestanddel at % PRINTED AS IS (Bestanddel: Bestanddel at: at for af (600 ppi))
hans »reale Selv«\footnote[2]{Se min „Psykologi“. Kap. V, 5; VI E, 2; VII B, 1 c.}. Og hos ethvert Individ vil der
tillige --- formedelst Erindring og Sammenligning, ja,
allerede formedelst det blotte Selvopholdelsesinstinkt ---
gøre sig en Sammenhæng gældende mellem de enkelte
% --- p. 9 ---
\setcounter{footnote}{0}%
\opage{9}Øjeblikke og de enkelte Livsperioder. Heri ytrer sig det
»formale Selv«\footnote[1]{„Psykologi“. Kap. V, 5 (cfr. IV, 7 e).}. Der vil derfor ytre sig en stærkere
eller svagere Tendents til at vurdere Handlingerne ikke
blot efter deres Indflydelse paa det enkelte Øjeblik eller
den enkelte Periode, men efter deres Indflydelse paa Livet
i dets Helhed. Og hos de forskellige Individer ytrer der
sig ikke blot Følelser, der ere knyttede til deres Bestaaen
hvert for sig alene, men ogsaa saadanne, som ere bestemte
ved de fælles Livsvilkaar, og som umiddelbart knytte det
ene Individ til et eller flere andre. Moderinstinktet og
de af det udspringende Følelser ere umiddelbare Naturtilskyndelser
i denne Retning. Inden for Familien, Stammen
og Nationen, og inden for de forskellige Associationer
og Kammeratskaber udvikle sig under Indflydelsen af
fælles Arbejde og fælles Skæbne mere eller mindre beslægtede
Følelser. Saadanne Følelsers Inderlighed og
Omfang staa i Regelen i omvendt Forhold til hinanden;
stærkest og kraftigst røre de sig i de snevreste Kredse,
især i Familien og Naboskabet. Men Historien viser, at
de have kunnet udvikle sig til at omfatte hele Menneskeheden.
Ideen om en Menneskehed og om almindelig
Menneskekærlighed (caritas generis humani) udviklede
sig allerede hos Oldtidens Grækere og Romere i de sidste
Aarhundreder før Kristi Fødsel, vistnok som Følge at % PRINTED AS IS (Kristi: som Følge at: at for af (600 ppi))
den Maade, hvorpaa de nationale Skranker sønderbrødes
 ved Alexander den Stores og Romernes Erobringer, saa
at Følelsen af Slægtskab og Fællesskab blev mulig.

De etiske Domme maa nu bestemmes ved en af de
her omtalte Følelser (foruden at de, som allerede anført,
bestemmes ved Tænkningens Konsekvents og Kendskabet
til de faktiske Forhold). Men deri ligger da, at \emph{der er
Mulighed for lige saa mange etiske Systemer},
% --- p. 10 ---
\setcounter{footnote}{0}%
\opage{10}\emph{som der er forskellige herskende Følelser, der
kunne bestemme Dommene}. De forskellige etiske
Systemer adskille sig fra hverandre ved Vurderingsmotivet.
Etikens Historie viser os, at de forskellige Muligheder
for en stor Del ere blevne benyttede. Endog
en ren Øjebliksetik, der lader Hensynet til det enkelte
Øjeblik, afset fra alle andre Øjeblikke, være eneraadende,
er bleven forsøgt (af Aristippos fra Kyrene). Man har
ofte taget Udgangspunktet i det enkelte, isolerede Individs
Selvopholdelse, Livsnydelse og Magtfølelse. Derved
føres man til et System, hvor Hensynet til andre Individer
kun faar Værdi som ofte uundværligt Middel til
at naa Goder for Individet selv. Etiken bliver altsaa en
Lære om de nødvendige Omveje, den Enkelte maa gaa
for at naa sin Lykke. Hvor Grundlag og Vurderingsmotiv
hentes fra Sympatien, kunne igen forskellige Systemer
blive mulige, alt efter det Omfang, Sympatien har.
En Familieetik haves tilnærmelsesvis realiseret i den
gammelindiske Etik. En national Etik er udviklet af
Platon og Aristoteles. En Stands- eller Associationsetik
fremtræder tilnærmelsesvis i Historien, ogsaa i vore Dage,
under de forskellige Samfundsklassers indbyrdes Brydninger.
Men ogsaa en Humanitetsetik arbejder sig frem
igennem Historien; Vurderingen af Handlinger og Institutioner
tager da deres Indflydelse paa alle menneskelige
Væsener i Betragtning.

\emph{Psykologisk og historisk set} er der en kontinuerlig
Sammenhæng mellem alle disse forskellige Standpunkter,
der i Teori og Praxis fremtræde gennemført
med mere eller mindre Konsekvents og med større eller
ringere Sagkundskab. Sammenhængen beror paa, at det
ene Vurderingsmotiv i Kraft af psykologiske Love og
under Indflydelse af Livets Erfaringer kan afløses af det
andet. Med en saadan Forandring forandres ogsaa Sy%
% --- p. 11 ---
\setcounter{footnote}{0}%
\opage{11}stemet af de etiske Domme; men denne Forandring er
Virkning, ikke Aarsag. Det er nye Følelsers Opstaaen,
det kommer an paa, og de opstaa ved praktiske Erfaringer,
ikke ved teoretiske Ræsonnementer. Der finder nu stadigt
en praktisk Opdragelse Sted, ikke blot inden for Familien,
men ogsaa inden for Samfundslivet i det Hele. Den
Enkelte lærer at underordne Hensynet til det enkelte
Øjeblik under Hensynet til Livshelheden, og Hensynet
til sin Selvopholdelse og Nydelse under Hensynet til,
hvad det fælles, almene Vel kræver. En Tilslutning til
andre Individer, der indlededes af egoistiske Motiver,
kan under det fortsatte Samliv og dets Erfaringer blive
selvstændig Genstand for Lystfølelse; hvad der først
grebes som Middel, kan blive Formaal. I den Art Motivforskydning
bestaar for en stor Del Menneskehedens Opdragelse
gennem Historien. Men den vilde neppe føre
vidt, naar den ikke støttedes ved umiddelbare sympatiske
Instinkter\footnote[1]{Smlgn. „Psykologi“, VI C, 2---7.}.

\emph{Etisk set} er der derimod Diskontinuitet mellem de
forskellige Standpunkter og Systemer, idet det ikke er
muligt for det ene at gendrive det andet ad den logiske
Bevisførelses Vej eller ved Hjælp af objektive Kendsgerninger.
Jo større Konsekvents det enkelte System besidder,
des uangribeligere vil det være. Den resolute og
konsekvente Egoist, der gør sig selv til eneste Formaal
og betragter andre Individer blot som Midler for sig selv,
kan ikke gendrives fra Familieetikens, Nationsetikens
eller Humanitetsetikens Standpunkt. Han siger med Shakesperes
Richard III: I am myself alone! Naturligvis er
det kun, for saa vidt Systemet er konsekvent, at det er
uangribeligt; hvis Familiefølelsen, Nationalfølelsen eller
Menneskekærligheden rører sig hos Egoisten, vil han
% --- p. 12 ---
\setcounter{footnote}{0}%
\opage{12}kunne komme i Modsigelse med sig selv og let gendrives.
Ellers kan han kun bringes til at opgive sit Standpunkt
ved, at der vækkes en Følelse hos ham, han hidtil ikke har
kendt. Men dette sker ved »Moralprædiken« (i højeste
Forstand), ikke ved blot »Moralbegrundelse«. Det er den
etiske Kunst, ikke den etiske Videnskab, som da træder
i Virksomhed.

Der gives i ethvert Samfund Individer, der ikke staa
paa det Grundlag, som er nødvendigt for, at Samfundslivet
paa det givne Trin kan bestaa. Saadanne Individer
ere Børn, Sindssyge og Forbrydere. Man kan sige, at
de have andre etiske Systemer end det, Samfundet som
saadant hylder. Over for dem anvendes der andre Midler
end det blotte Ræsonnement. De blive Genstand for en
praktisk Behandlingsmaade, en Tugt og Opdragelse, der
skal vække de Følelser hos dem, de hidtil manglede.
Samfundet tiltager sig i \emph{Kraft af sin egen Etik}, Ret
til at behandle dem paa denne Maade, men kan ikke
altid begrunde denne Ret ud fra saadanne Individers egen
Etik. --- I enhver etisk Dom, vi udtale, er der egentligt
et saadant pædagogisk Moment. Den etiske Dom udspringer
ikke blot af Følelsen; den taler ogsaa til Følelsen.
Den, hvis Handling vi vurdere, rose eller dadle vi;
vi sætte derved hans Følelse i Bevægelse, gøre os til
Autoritet over for ham og have Ansvaret for Virkningen
af de Motiver, vi paa denne Maade kunne fremkalde.

4. Imidlertid maa det vel mærkes, at de forskellige
Systemer ikke ere hverandre rent sideordnede. Det enkelte
System kan give de andre Systemer en vis underordnet
Betydning inden for sig. Egoistens System f. Ex. udelukker
ikke Hensynet til Familie og Stat, men betragter
dem blot under Synspunktet af Midler. Nations- og Humanitetsetiken
maa fra deres Standpunkter vise, hvorledes
det enkelte Individ inden for deres Systemer kommer til
% --- p. 13 ---
\setcounter{footnote}{0}%
\opage{13}sin fulde Ret, skønt paa anden Vis end efter det egoistiske
System. Gennem Sympatien, gennem Hengivelsen
til et større Hele og Arbejde i dets Tjeneste udvikles
den Enkeltes Kræfter og Evner; hans Personlighed udvides
og styrkes netop gennem Selvforglemmelsen. At
glemme sig selv er ikke det Samme som at tabe sig selv.
Det er at sætte sig andre Formaal end sit eget Jeg.
Ved den forøgede og den forhøjede Stræben, disse videre
gaaende Formaal betinge, faar Sindet større Liv, Kraft
og Flugt. Det er saaledes en Berigelse, en Vinding for
Individet, hvis Sympatien faar Magt i det. Men dette
kan først erfares, naar Sympatien \emph{har afløst} Egoismen,
og kan altsaa ikke være Grund eller Motiv for denne
Afløsning. Derfor vilde Sympatien ikke have udviklet
sig i Menneskeslægten, dersom Naturen ikke her, lige
som ved Selvopholdelsen, greb fat i os gennem Instinkter
og uvilkaarlige Erfaringer. Vi vilde ellers i al Evighed
staa som Manden, der ikke vilde gaa i Vandet, før han
kunde svømme.

Man kunde maaske mene, at det her gik som ved
et fysisk eller kemisk Experiment, ved hvilket der ad
Erfaringens Vej (a posteriori) godtgøres Noget, som man
ad den blotte Tænknings Vej (a priori) ikke kunde godtgøre,
og at Etiken altsaa her ikke var stillet anderledes
end enhver Videnskab, der optager Forudsætninger fra Erfaringen.
Men man vilde da overse en meget væsentlig
Forskel. Det Experiment, der her er Tale om, skal foregaa
med Individet selv, og netop med det i ham, der
betinger hans etiske Dom. Naar Individet har erfaret,
hvorledes Sympatien udvider hans Jeg, gør det sundere,
kraftigere og aabner det nye Kilder til Tilfredsstillelse,
saa er han selv bleven en Anden end han før var; han
har faaet et andet realt Selv. Han gør Regnskabet op
fra et andet Grundlag end det, han stod paa, før han
% --- p. 14 ---
\setcounter{footnote}{0}%
\opage{14}foretog Experimentet, --- og Opgørelsen er helt igennem
bestemt ved det nye Grundlag, det nye Vurderingsmotiv.
Han ændrer altsaa ikke sin Anskuelse ved Slutning fra
Experimentet, men dette ændrer umiddelbart hans Følelse
og som Følge deraf hans Anskuelse.

Desuden sker Overgangen fra det ene Grundlag til
det andet ikke med samme Lethed som den, med hvilken
et Experiment foretages. Det viser Historien og daglig
Erfaring tilstrækkeligt. Menneskehedens Opdragelse tager
lang Tid. For nogle Individers Vedkommende har man
endog benegtet Muligheden af, at der kan vækkes det
Minimum af Sympati, det nu bestaaende borgerlige Samfund
kræver; i en anden Sammenhæng ville vi drøfte
denne Benegtelse nærmere. Erfaringen viser i hvert Tilfælde,
at der gives Individer, som i Konsekvents af det
Grundlag, hvorpaa de nu én Gang staa, tørne sammen
med de givne Forhold og gaa til Grunde. Tragedien
henter nogle af sine største Emner fra etiske Systemers
Sammenstød; inden for et og samme System er tragisk
Sammenstød ikke saa let muligt. --- Og selv hvor Overgangen
fra et Grundlag til et andet lykkes, tager den i
Regelen Tid og sker ofte gennem stærke Kriser.

Det vilde derfor være en meget uhensigtsmæssig Brug
af Ordene, hvis man vilde kalde en Følelse egoistisk,
fordi det dog ogsaa i Sympatien er vor \emph{egen} Lyst eller
Ulyst, vi have, og fordi Sympatien ikke kunde bestaa,
naar den ikke medførte et fyldigere og inderligere Liv
og derved en større Tilfredsstillelse og Selvhævdelse end
der vilde opnaas, naar Lyst- og Ulystfølelse kun bestemtes
ved en Selvopholdelse og Selvhævdelse, for hvilken alt
Andet i Verden var blotte Midler. En saadan Brug af
Ordene vilde forlede til at overse nogle af de betydeligste
og dybest indgribende Modsætninger paa det aandelige
Livs Omraade.

% --- p. 15 ---
\setcounter{footnote}{0}%
\opage{15}5. Det beror paa et Valg, hvilket Standpunkt der
indtages i en Fremstilling af Etiken. Det vilde have sin
store Interesse at faa de forskellige mulige Systemer hvert
for sig fuldt udførte, saa at vi kunde se, hvorledes Etiken
former sig i det Enkelte paa et Øjebliksstandpunkt, et
egoistisk Standpunkt, et Familiestandpunkt o. s. v. Det
vilde kaste Lys over Vurderingsmotivernes Natur. Den
historiske eller sammenlignende Etiks Interesse beror
herpaa; ti historisk finde vi Tilnærmelser til de forskellige
mulige Systemer. Men ogsaa kun Tilnærmelser. Ti i det
virkelige Liv ere Vurderingsmotiverne (lige saa vel som
Handlingsmotiverne) altid mere eller mindre sammensatte.
Vor Samvittighed (det vil sige, den Følelse, der faktisk
bestemmer vor Vurdering) falder neppe nogensinde fuldstændigt
sammen med Sympatien eller Egoismen eller med
en enkelt af deres Former. Følelser af højst forskellig Art
virke sammen i den. Naar Etiken for at naa fuld Klarhed
og Konsekvents udtager en enkelt bestemt Følelse
og holder sig til den som Grundlag, idet den tænker sig
den i fuld Renhed og Styrke og med største Omfang,
foretager den en lignende Abstraktion og Idealisation,
som enhver formel Videnskab foretager ved sit første
Skridt. Det rene Rum, som Geometrien forudsætter i sine
første Definitioner og Axiomer, forekommer ikke i Virkeligheden
som Genstand for Sansning og Anskuelse.

Naar nu Etiken til Grundlag vælger en ideal Samvittighed,
der i sin Vurdering bestemmes ved universel
Sympati, har den vel foretaget en Abstraktion, men denne
bestaar i, at den tænker sig en Kraft i ideal Fuldendelse,
som faktisk bestaar i Verden, har arbejdet sig frem og
sandsynligvis vil vedblive at arbejde sig frem. Den har
altsaa trods sin Idealitet fast Grund under Fødderne.
Den foreslaar os at betragte den menneskelige Verden
og det, som dér foregaar, fra dette Stade og se, hvilke
% --- p. 16 ---
\setcounter{footnote}{0}%
\opage{16}Idealer og Opgaver deraf følge. Saa kunne Andre jo betragte Verden fra deres Stade. Naar vi da hver for sig
gaa frem med saa stor Konsekvents som muligt og støtte
os til saa mange og sikre Erfaringer som muligt, vil vor
Erkendelse have gjort de Fremskridt, den paa dette Omraade
formaar at gøre. Muligheden at at begrunde Etiken % PRINTED AS IS (begrunde: Muligheden at at begrunde: for af at (600 ppi))
fører os til Nødvendigheden af at anerkende en Begrænsning
for dens Videnskabelighed. Begrundelse og Begrænsning
høre paa alle Omraader nøje sammen. ---

Naar Sympatien vælges som Vurderingsmotiv,
ligger deri, at kun de Formaal, som en oplyst, konsekvent
og universel Sympati vedkender sig, kunne være
de højeste etiske Formaal. Men der ligger ikke deri, at
Sympatien nødvendigvis bør være Handlingsmotiv i
alle Tilfælde og hos alle Individer. Der kan i den menneskelige
Natur ligge andre Motiver, der ikke behøve at
føre til Modstrid mod de Formaal, Sympatien sætter, men
netop kunne arbejde til Bedste for dem; og disse Motiver
kunne endog være uundværlige. Trang til Selvopholdelse
og Selvhævdelse, Ærgerrighed, Tilslutning til Sæd og
og Skik eller til ærværdige Autoriteter --- det er Exempler
paa Kræfter og Tendentser i den menneskelige
Natur, som en Etik, der støtter sig til Sympatien som
sidste Vurderingsmotiv, fuldt ud kan anerkende, for saa
vidt de virke i Retning af de Formaal, den har sat.
Forskellen mellem saadanne Motiver og den universelle
Sympati er, at medens denne sidste \emph{nødvendigvis}
kommer i Modsigelse med sig selv ved ikke at betragte
Handlinger og Institutioner fra Helhedens (Samfundets,
tilsidst Menneskehedens) Synspunkt, komme de andre Motiver
ikke nødvendigvis i Modsigelse med sig selv ved
en mere begrænset Synskreds. For Egoisten er det ligegyldigt,
hvorledes det gaar med de Andre (hvad enten
de ogsaa ere Egoister eller ikke), naar de blot fremme
% --- p. 17 ---
\setcounter{footnote}{0}%
\opage{17}\emph{hans} Formaal. Og selv om der, som Nogle have ment
at kunne godtgøre, skulde være en Harmoni mellem alle
Egoisters velforstaaede Interesser, saa glæder den enkelte
Egoist sig kun herover, for saa vidt denne Formaalenes
Harmoni kommer \emph{ham} til Gode. Dersom Forestillingen
om en saadan virkelig eller mulig Interessernes Harmoni
hos ham vakte en Glæde eller maaske endog Begejstring,
som gaar ud over den Glæde, der føles over et Middel,
der kun har Værd som Middel, saa vilde han ikke mere
være Egoist; han vilde da være naat til at have andre
Formaal foruden sig selv.

Den paa Sympatien som Vurderingsmotiv grundede
Etik fører hver Enkelt til at betragte sig ikke blot som
Formaal, men ogsaa som Middel, og til at betragte Andre
ikke blot som Midler, men ogsaa som Formaal. Fra intet
andet Grundlag kan denne Betragtningsmaade begrundes.
Dens almindelige Grundtanke kan udtrykkes saaledes:

Se paa dit Liv og paa de Andres Liv ikke blot fra
dit eget Stade, som om det var Verdens Midtpunkt, men
fra et højere Stade: Samfundets, Slægtens! Fordi du
naar det Maal, du har sat dig, er det ikke sagt, at alle
Andre ogsaa have naat deres Maal. Fordi du har faaet
dine Ønsker opfyldte, er der Lidelse nok endnu tilbage i
Verden, berettigede, men uopfyldte Ønsker i tilstrækkelig
Mængde. Var du end den Største, der har levet, saa er
du dog kun én blandt Mange af den Slægt, hvis Historie
ikke er afsluttet med din Historie!

6. En væsentlig Indvending mod Forsøget paa at
bygge Etiken paa Sympatien som Vurderingsmotiv, kunde
synes at ligge i Kampen for Tilværelsen som ledende
Synspunkt for alt Liv, ogsaa det menneskelige Liv. Naar
Livet er en uafbrudt Kamp for at bestaa mellem forskellige
Arter og mellem forskellige Individer af samme
Art, gælder det saa ikke for hver Enkelt at komme de
% --- p. 18 ---
\setcounter{footnote}{0}%
\opage{18}Andre i Forkøbet og at skaffe sig saa stor Fordel som
muligt paa deres Bekostning, og maa en paa Sympati
grundet Etik saa ikke savne al Rod i det virkelige Liv?
Der synes jo dog saa ikke at gælde anden Modsætning
end den, der bestaar mellem dem, der kunne bestaa og
udvikle sig, og dem, der ikke kunne det, --- mellem de
Sunde, Kraftige og Lykkelige paa den ene Side, de Syge,
 Svage og Lidende paa den anden Side. Og er der her
ikke en absolut, en haabløs Modsætning, ved hvilken der
ikke kan rokkes, uden at almindelig Svækkelse og Lammelse
bliver Følgen?

Medens Nogle fortolke den Darwin'ske Lære om
Kampen for Tilværelsen paa denne Maade og med æstetisk
Velbehag svælge i Tanken om den Hensynsløshed
og Kraft, Tilværelseskampen fører til at udfolde, er der
Andre, som vel indrømme, at saadanne etiske Konsekventser
ligge i hin Lære, men som deri finde en af de
væsentligste Indvendinger imod den. --- Hvad Darwin
selv angaar, indrømmede han ingen af disse Konsekventser.
Som Moralfilosof sluttede han sig til den af Shaftesbury
og Hutcheson grundlagte og i vore Dage af Herbert
Spencer videre førte etiske Skole, der i Sympatien saa
Etikens Grundlag, en Opfattelse, som det i denne Afhandling er forsøgt at give en bestemtere Belysning og
Gennemførelse. Jeg skal nu forsøge i Korthed at vise,
at den store Naturforsker har havt en rigtigere Opfattelse
af sin Lære end den, der kommer frem hos dem,
der drage hine formentlige Konsekventser.

7. Kampen for Tilværelsen bringer alle Muligheder
til Udfoldelse, som ere til Hjælp under de givne Forhold.
Der kan saaledes udvikle sig en hensynsløs Selvopholdelsestrang,
der kan stige til forskellige Grader af bevidst
Egoisme. Der ligger en Mulighed og en Tendents i den
menneskelige Natur i denne Retning. Men det er ikke
% --- p. 19 ---
\setcounter{footnote}{0}%
\opage{19}den eneste Mulighed, som ligger i den menneskelige
Natur, og heller ikke den eneste Mulighed, som kan udfolde
sig under Kampen for Tilværelsen. Denne Kamp
er ikke i bogstavelig Forstand en »Kamp af Alle mod
Alle«, saaledes som Thomas Hobbes i sin Tid antog.
Den Tilværelse, som Kampen gælder, er ikke blot det
enkelte Individs, men Artens, Slægtens. Og Arten opholdes
ikke altid og ikke blot ved, at enkelte Individer
hvert for sig saa hensynsløst som muligt søge at holde
sig oppe, men ofte ogsaa, ja: kun, ved at de stræbe og
kæmpe i Fællesskab, eller ved, at det ene Individ ofrer
sig for det andet. Kampen for Tilværelsen frembringer
altsaa ikke blot Egoisme, men ogsaa Midler til at komme
ud over Egoismen. Dens Virkninger ere ikke alle enstydige,
ikke alle af samme Værdi i etisk Henseende. En
Etik er netop nødvendig, fordi der maa træffes et Valg
mellem disse Virkninger, saa snart Spørgsmaalet er om,
hvilke af dem vi med vor bevidste Villen og Indgriben
skulle begunstige, og hvilke vi skulle søge at hæmme
eller fortrænge. Som villende Væsener ere vi ikke blot
Virkninger af den forudgaaende Udvikling, men ogsaa
(for en Del i det Mindste) Aarsager til den følgende Udvikling,
og vore etiske Overvejelser angaa den Retning,
i hvilken vor Kausalitet bør gaa. For at afgøre dette,
behøve vi en Maalestok, et Ideal, og Etikens første Opgave
er at fastsætte det Grundlag, ud fra hvilket en
saadan Maalestok kan findes. Den filosofiske Etik søger,
som vi have set, dette Grundlag i den menneskelige
Natur. Kampen for Tilværelsen kan, saa sandt den er
Betingelsen for alt Liv, ikke stanses; men den kan
maaske føres paa en Maade, der stemmer med det givne
Vurderingsmotiv.

Ved at se paa de første Betingelser for Slægtens
Bevarelse ville vi finde en Mulighed herfor.  --- Goethe
% --- p. 20 ---
\setcounter{footnote}{0}%
\opage{20}sagde, da han saa et Billedhuggerarbejde, der fremstillede
en Ko med en diende Kalv: »Her have vi et Emne af
den højeste Art: det Verden opholdende, gennem hele
Naturen gaaende ernærende Princip træder os her i en skøn
Lignelse i Møde«. Den Darwin'ske Lære er en nærmere
Udførelse af denne Tanke. Den har paavist, at uden Instinkter,
der føre til Omsorg og Opofrelse for Afkommet,
kan en Art ikke bestaa i Kampen for Tilværelsen. Saa
mange Gaader, der end staa tilbage, hvad de første Spirer
til det vidunderligste Instinkt i Naturen, Moderinstinktet
eller Forældreinstinktet, angaar, saa er dets Nødvendighed
indlysende. Det viser sig især, hvor det staar i den grelleste
Modsætning til et Væsens andre Instinkter. Edderkopperne
f. Ex. ere Rovdyr, der søge deres Bytte hver
for sig ved Hjælp af deres Net, uden at der, som hos
saa mange andre Rovdyr, kan udvikle sig en vis social
Følelse under Jagtkammeratskabet. De angribe og fortære
tvertimod andre Individer af deres egen Art, selv
om der ikke er Mangel paa andet Bytte. End ikke i
Parringstiden ere Han og Hun sikre for hinandens Rovlyst.
Afkommet vilde derfor strax gaa til Grunde, dersom
ikke selv hos disse ellers saa »egoistiske« Dyr Moderinstinktet
fremtraadte endog i usædvanlig høj Grad\footnote[1]{Se \emph{Fr. Dahl}: Versuch einer Darstellung der psychischen Vorgänge in den Spinnen. (Vierterljahrsschrift für wissenschaftliche Philosophie. IX). p. 180 ff.}. ---
De Instinkter og Følelser, som knytte Forældre og Børn
til hinanden, begunstiges altsaa afgjort ved Kampen for
Tilværelsen. Kun i sjeldnere Tilfælde ere de modsatte Instinkter
Genstand for naturlig Udvælgelse. Saaledes kan
det Instinkt, der faar Arbejdsbierne til at dræbe deres
Brødre, Dronerne, og Bidronningen til at dræbe Biprin%
% --- p. 21 ---
\setcounter{footnote}{0}%
\opage{21}sesserne\footnote[1]{\emph{Darwin}: Menneskets Afstamning. I. p. 75. --- Arternes Oprindelse. p. 259.}, kun forklares ved dets Betydning under Tilværelseskampen,
der altsaa under nogle Omstændigheder
begunstiger Moderkærlighed og Søskendekærlighed, under
andre Omstændigheder det Modsatte; dog er dette Sidste
det sjeldneste Tilfælde. Hvor inden for en og samme Art
enkelte Individer i større eller mindre Grad mangle
Forældreinstinktet,\footnote[2]{Se herom \emph{Lombroso}: Der Verbrecher. Deutsche Uebers. p. 9. 27. --- Dictionnaire des sciences anthropologiques. Article „Famille“.} maa saadanne Undtagelser forklares
paa samme Maade, saa fremt de ikke ere Rester fra tidligere
Udviklingstrin, der holde sig i enkelte Tilfælde
formedelst Atavisme.

Denne Betragtningsmaade finder ogsaa Anvendelse
paa Sammenholdets og Fællesskabets Betydning for de
enkelte Individer under deres Livsopholdelse og Livsførelse.
Ikke som Følge af bevidst Beregning og vilkaarligt
Valg, men i Følge Forholdenes Natur kommer
det enkelte Individ til at føre Kampen for Tilværelsen
sammen med andre. Det indøver praktisk Sammenholdet
og Samarbejdet, før det kan danne sig nogen Forestilling
om dets Betydning. Der danner sig herved inden for en
og samme Gruppe af Mennesker en Broderskabs- og
Menneskelighedsfølelse, som maaske fra først af kun
viser sig under Krig og Fare, men senere kan ytre sig
ogsaa under andre Forhold.

Under det saaledes uvilkaarligt indstiftede og indøvede
fælles Liv, under Fællesskabet i Skæbne og Arbejde,
i Lyst og Smerte, lærer den Enkelte praktisk at
anvende en højere Maalestok end den, der kan hentes fra
hans egne, rent individuelle Interesser. Gennem Samlivet
og Samvirksomheden udvikler Samfølelsen sig, --- lige%
% --- p. 22 ---
\setcounter{footnote}{0}%
\opage{22}gyldigt hvilke Motiver, der fra først af have ført til
Samlivet og Samvirksomheden, ti en Motivforskydning
finder her ofte Sted.

I min »Etik« har jeg vist, hvorledes denne praktiske
Udvikling af Fællesskabsfølelsen spiller en Rolle inden
for Ægteskabet (p. 180), ved Arbejderstandens Organisation
(p. 244), ved Religionsfrihedens Udvikling (p. 312)
og ved Nationalfølelsen (p. 338 f.). Jeg har ført den tilbage
til »det aristoteliske Princip« (p. 162 f.), at uvilkaarlig
Indøvelse af Egenskaber og Færdigheder gaar
forud for bevidst Besiddelse og Udøvelse. Egenskaber
danne sig, --- det var Aristoteles's, ved den nyere Biologi
og Psykologi bekræftede Lære, --- ved Virken og
Handlen, og hvor en Egenskab eller en Færdighed har
dannet sig, er der forbunden en Lystfølelse med Udøvelsen
af de tilsvarende Handlinger. Den første Betingelse
er, at Livsforholdene giver Anledning til Øvelse.

8. Endnu et Hovedpunkt er der at fremhæve til
Belysning af Etikens Forhold til Kampen for Tilværelsen.
Denne gælder ikke blot den rene, nøgne Existents, den
blotte fysiske Bestaaen. Den Tilværelse, der kæmpes om,
er den, som føres paa et bestemt Livstrin. Ikke uden
Modstand, ikke uden at flere Muligheder først prøves,
gaar Livet tilbage til et lavere, mere elementært Trin.
Ethvert Livstrin har sine Fornødenheder, i hvilke dets
Væsen lægger sig for Dagen. Kun paa de aller simpleste
Livstrin er Lyst- og Ulystfølelsen ene knyttet til, hvad
der umiddelbart tjener til at opholde Livet. Naar disse
mest elementære Fornødenheder ere tilfredsstillede, vil
der paa højere Trin røre sig en Trang til at bruge
Kræfterne og en Lyst ved denne Brug, selv om den ikke
fører til noget praktisk Resultat. Lystfølelsen ved Leg
og den æstetiske Følelses Udvikling ere Exempler herpaa.
Glæden ved Sandhedserkendelse er ogsaa et Exempel paa
% --- p. 23 ---
\setcounter{footnote}{0}%
\opage{23}denne Motivforskydning, ved hvilken Lystfølelsen knyttes
til selve Virksomheden og ikke blot til dennes Resultat.
Der er naat et højere Trin, hvor Livets elementære Krav
træde saaledes tilbage, at ideelle, af det fysiske Selvopholdelseskrav
uafhængige Interesser blive mulige. Og
naar nu et saadant Trin er naat, gælder Kampen for
Tilværelsen dette Livstrins Bevarelse: Livet er det, der
føres paa dette Trin; alt Andet er Død. Med Rette taler
man derfor om en aandelig Død, hvor ideelle Interesser
ere udslukte. Til saadanne ideelle Interesser høre nu
ogsaa Sympatien og Samvittigheden, naar de have arbejdet
sig frem. Et Liv, som kun vindes ved at fornegte
dem, vilde nu staa som værdiløst\footnote[1]{Smlgn. „Om Grundlaget for den humane Etik“. (København 1876), p. 30 f., hvor jeg allerede har fremsat denne Tanke.}.

% --- p. 24 ---
\setcounter{footnote}{0}%
\essay[Om Velfærdsprincipet]{Om Velfærdsprincipet.}
\opage{24}


1. Naar Velfærdsprincipet anerkendes som \emph{Maalestok}
ved Vurderingen af Handlinger og Institutioner,
vil det sige, at disse kaldes gode eller onde, alt efter
som de i deres Virkninger (saa langt vi kunne forfølge
disse) medføre overvejende Lystfølelser eller overvejende
Ulystfølelser hos en større eller mindre Kreds af bevidste
Væsener. Enhver Handling kan sammenlignes med en
Sten, som kastes i Vandet; Bevægelsen forplanter sig i
større eller mindre Kredse, og Vurderingen beror paa,
hvorledes den virker dér, hvor den griber ind i levende
Væseners Ve og Vel. Lige som den teoretiske Videnskab
forklarer det enkelte Naturfænomen ved dets Sammenhæng
med andre Naturfænomener, saaledes vurderer Etiken den
ene Følelse ved dens Forhold til andre Følelser: den
Tilfredsstillelse, den Handlende selv føler ved sin Handlen,
kan kun billiges, dersom Handlingens Virkninger ikke
gribe forstyrrende ind i andre Væseners Lystfølelse, naar
et saadant Indgreb da ikke kan paavises som nødvendigt
Middel for en saa meget des større Lystfølelse, --- og de
bedste Handlinger ere de, som forøge Lystfølelsen i Verden
og formindske Ulystfølelsen. --- Dette Princip svarer som
\emph{Maalestok} umiddelbart til Sympatien som \emph{Vurderingsmotiv}.
Hvor vidt det er muligt at anerkende det
% --- p. 25 ---
\setcounter{footnote}{0}%
\opage{25}fra Standpunkter, der lægge andre Vurderingsmotiver
til Grund, skal her ikke nærmere drøftes\footnote[1]{Se herom „Etik“ p. 39---42.}.

Vurderingen bliver ikke staaende ved den ydre Handling,
men gaar tilbage til den Handlendes Motiver, til
hans Karakteregenskaber, til hele det indre Liv, af
hvilket Handlingen er udsprungen. Dette har sin Grund
i dens Natur og Oprindelse. De etiske Domme ere fra
først af uvilkaarlige Følelsesudbrud. Og disse Følelsesudbrud
ere af stor praktisk Betydning, idet de virke
bestemmende paa Villien (den egne og Andres) og frembringe
Motiver til fremtidige Handlinger. Det er da nødvendigt
at gaa tilbage til det Punkt, hvor det er muligt at
paavirke den Kraft, som har frembragt Handlingen, og
dette Punkt ligger i det indre Liv, i Sindelaget hos den
Handlende, i de Motiver, der have hidført Handlingen.
Kun naar de ændres, ændres ogsaa de fremtidige Handlinger,
og med disse igjen deres Virkninger. Det er
Virkningstendentsen, der konsekvent forfølges tilbage til
sin Kilde.

En Kraft kende vi kun fra dens Virkninger. Kun
ved Slutning kunne vi danne os en Forestilling om hvad
der foregaar i et Menneskes Sind, om de Motiver og
Evner, han raader over. En Godhed eller Storhed, der
aldrig nogensinde viste sig i Handling, vilde ikke kunne
blive Genstand for etisk Anerkendelse. Den vilde bero
paa Selvbedrag. Der maa dog i det Mindste foregaa et
indre Arbeide, en Higen og Stræben i den Retning, det
etiske Princip kræver. Ofte maa Individet i sin Vurdering
af sig selv tage sin Tilflugt til dette indre Arbejde,
og den sande etiske Vurdering maa tilsidst angaa det\footnote[2]{Se herom den følgende Afhandling om Forholdsloven i Etiken.};
men det er jo ogsaa netop Begyndelse, Tilløb til det, som
% --- p. 26 ---
\setcounter{footnote}{0}%
\opage{26}Velfærdsprincipet fordrer, kun at Fortsættelsen ofte paa
Grund af individuelle Forhold umuliggøres.

Men lige saa vigtigt er det at fastholde, at Virkningen
af Handlingen ikke strax behøver at vise sig.
Naar store og gode Mennesker saa ofte miskendes af
deres egen Samtid, forklares det ved, at kun et meget
omfattende Blik kan udmaale Betydningen af deres Færd.
Deres Storhed og Godhed beror netop paa, at deres Tanke,
Følelse og Villie omfatter langt mere end deres kortsynede
og snevertfølende Samtid formaar at opdage. Lang
Tid kan gaa, inden almindelig Forstaaelse og Tilegnelse
bliver mulig. Derfor ligger det ingenlunde i Velfærdsprincipet,
at man i sin Færd skal rette sig ganske efter
den Trang og de Fornødenheder, Menneskene i Øjeblikket
have. Velfærdsprincipet fordrer netop, at vi ikke sky
Striden mod Fordomme og Træghed. Det Bedste, man
kan gøre for Andre, er ofte simpelt hen at lade dem føle,
at de endnu i deres Ønsker og Fornødenheder staa for
lavt og ikke stille tilstrækkeligt store Fordringer. Saaledes
gaar --- for at tage et Exempel --- den store Kunstner
ofte sin ensomme Gang, uforstaaet eller endog misforstaaet
af den store Mængde. Og dog følger han Velfærdsprincipet,
uden at han maaske tænker derpaa, ved strængt
at hævde Kunstens Fordringer. Han forøger Slægtens
aandelige Kapital og skænker den en Kraft, som i Fremtiden
kan komme til at virke i store Kredse. Kun en
kortsynet Opfattelse og Anvendelse holder sig til den i
Øjeblikket herskende Trang i Stedet for at se hen til
Slægtens varige Livsbetingelser og til de stadige Kilder
for nyt Liv og ny Virksomhed\footnote[1]{Denne sidste Udvikling er her hidsat fra min „Etik“ p. 94.}.

2. Velfærdsprincipet afgiver en \emph{Regel}, vi kunne
lægge til Grund ved vor Vurdering af Handlinger. Men
% --- p. 27 ---
\setcounter{footnote}{0}%
\opage{27}det fordrer ingenlunde, som man under Tiden har ment,
at Hensynet til Velfærd ogsaa skal være \emph{Motiv} for enhver
Handling. Vi gribe kun til almindelige Principer
for at orientere os i saadanne Tilfælde, hvor det umiddelbare
Skøn eller den instinktmæssige Følelse ikke kan bringe
Sagen til Afgørelse, altsaa i Tvivlstilfælde, eller ogsaa
naar vi tilstræbe en systematisk Behandling af etiske
Spørgsmaal. Den etiske Følelse kan virke aldeles uvilkaarligt
og uden egentligt Ræsonnement, idet vi umiddelbart
betages af Handlingen som den staar for os (for
Tanken, altsaa som mulig, eller i Erfaringen, altsaa som
virkelig), lige som vi i vor æstetiske Følelse umiddelbart
kunne gribes af et Kunstværks eller et Landskabs Skønhed,
uden at behøve at ty til æstetiske Ræsonnementer. Eller
man følger trygt de »uskrevne Love«, som indeholdes i
Sæd og Skik, i Tradition, i det Hele i den positive
Moralitet. Netop i Kraft af Velfærdsprincipet vil en
saadan Umiddelbarhed være at anbefale og opretholde, saa
længe den ikke fører til at overse virkelige Problemer
og Opgaver. Det er Uskyldighedstilstanden, af hvilken
Ingen bør udrives uden Nødvendighed. Naar den umiddelbare
Tryghed er brudt, blive de abstrakte Principer
nødvendige Hjælpemidler; men de kunne ofte kun gennem
mange Mellemled anvendes i de enkelte, konkrete Tilfælde,
og de faa ikke saa let en praktisk Indflydelse
paa Villien.

Endog helt andre Motiver end etisk Følelse eller
Hengivelse til den positive Moralitets Fordringer kan
Velfærdsprincipet billige, ønske eller kræve. Det er saaledes
heldigst, om en Mand drager Omsorg for sin Hustru
og sine Børn, fordi han elsker dem, og ikke fordi hans
etiske Følelse kræver det. Hvor Pligten bevidst skal
træde i Kraft i de inderligste Forhold mellem Mennesker,
er det i Regelen et Tegn paa, at et mindre heldigt Forhold
% --- p. 28 ---
\setcounter{footnote}{0}%
\opage{28}er til Stede. Den fuldkomne Kærlighed driver ikke blot
Frygten, men ogsaa Pligten ud. Lige som Familieforholdet
er ogsaa Kunstnerens og Forskerens Forhold til deres
Fag et Exempel herpaa\footnote[1]{Lige som der her er skelnet mellem Maalestok eller Vurderingsprincip og Handlingsmotiv, saaledes er der i den foregaaende Afhandling (§ 5) skelnet mellem Vurderingsmotiv og Handlingsmotiv. Paa Grund af den nøje Sammenhæng mellem Vurderingsprincip (Maalestok) og Vurderingsmotiv (Grundlag) falde disse to Betragtninger omtrent sammen.}.

Der er af Wundt\footnote[2]{Ethik. Eine Untersuchung der Thatsachen und Gesetze des sittlichen Lebens. Stuttgart 1886. p. 339.} opstillet følgende Indvending
mod Velfærdsprincipet: »Det er tænkeligt, at En opofrer
sig for en Anden; det er tænkeligt, at En hengiver Gods
og Liv for et eller andet idealt Formaal, for Fædrelandet,
Friheden, Troen eller Videnskaben. Men det er aldrig
forekommet og vil aldrig forekomme, at Nogen resignerer,
for at den Sum af Lykke, der gives i Verden, kan blive
større«. Denne Indvending overser netop, at hvad der
er Princip for Vurderingen af en Handling, naar denne
anerkendes som god, ikke \emph{behøver} at være Motiv for
Handlingen. Den Handlendes Tanke og Følelse kunne
meget godt stanse ved Fædrelandet, Friheden, eller hvilket
andet idealt Formaal han ofrer sig for, uden at han behøver
at anstille nogen Reflexion over hvad der betinger
selve disse ideale Formaals Værdi. Men i en systematisk
Etik eller i praktiske Tvivlstilfælde drøfte vi Fædrelandskærlighedens\footnote[3]{Se min „Etik“. Kap. XXXVI.},
Frihedsfølelsens\footnote[4]{Ibid. Kap. XXIII og XL.}, Troens\footnote[5]{Ibid. Kap. XXXI---XXXIII.} og Videnskabens\footnote[6]{Ibid. Kap. XXVIII og XXIX.}
Værd og Betydning for det menneskelige Liv.
Dersom f. Ex. Frihed ikke var et Gode for et Folk, vilde
% --- p. 29 ---
\setcounter{footnote}{0}%
\opage{29}det være urigtigt, om et Menneske vilde ofre Livet for
at bevare den. Naturligvis drejer det sig i det enkelte
Tilfælde aldrig om det abstrakte Begreb Velfærd i og
for sig, lige saa lidt som det ved en speciel teoretisk
Undersøgelse drejer sig om det abstrakte Begreb Aarsag
i og for sig. Men vi opstille i Etiken Velfærdsprincipet
og i Erkendelsesteorien Aarsagsprincipet, naar vi ad
Analysens Vej have søgt de første Forudsætninger for
vor praktiske og teoretiske (vurderende og forklarende)
Aandsvirksomhed. ---

3. Mod Velfærdsprincipet strider det ikke, at Lyst
saa ofte maa købes med Smerte. Ti Smerten staar da
kun som det nødvendige Gennemgangsled, og Velfærdsprincipets
Betydning er netop at hævde, at Bevisbyrden
paahviler den, der paastaar, at et saadant Gennemgangsled
er nødvendigt. Al Tilføjelse af Smerte skal motiveres,
medens Lystfølelsen i og for sig (naar ikke det, der fremkalder
den, tillige volder større uheldbringende Virkninger)
er berettiget. Velfærdsprincipet siger blot: frembring saa
megen Lyst og saa liden Ulyst som muligt ved din
Handlen! I hvilken Grad det er muligt at fyldestgøre
denne Fordring, derom indeholder Velfærdsprincipet i og
for sig Intet. Et Princip omstødes dog ikke ved Vanskeligheden
af at anvende det. Og Erfaring viser, at der gives
en Lystfølelse, som ikke er for dyrt købt, selv om der
gaar dyb Smerte med i Købet. Man har mod Velfærdsprincipet
anført Clärchens Sang (i Goethe's »Egmont«):
»Himmelhoch jauchzend, \emph{zum Tode betrübt}«. Men
lad Clärchen tale ud, og lad os lægge Mærke til Sangens
sidste Linier, i hvilke hun uddrager Resultatet af hele
sin Stemningsrække: »Glücklich allein ist die Seele, die
liebt!« Sagen er: der gives en Bevægelse i Sindet, et
Liv i Følelsen, som er forbundet med en saa dyb og stor
Tilfredsstillelse, at den stærke Svingning mellem Lyst og
% --- p. 30 ---
\setcounter{footnote}{0}%
\opage{30}Smerte ikke ophæver den totale Lykkefølelse, men netop
styrker den. Der virke her to psykologiske Love sammen.
Den ene er den, i Følge hvilken Smerten (eller Ulysten,
eller Sorgen), naar den ikke overstiger en vis Grænse,
kan danne Baggrund for Lystfølelsen og derved forhøje
denne. Allerede heri ligger et tilstrækkeligt Motiv til at
vælge saadanne Tilstande i Stedet for dem, der ikke naa
saa højt i Intensitet, skønt de paa den anden Side ere
mere usammensatte Lysttilstande. Den anden Lov er den,
at der i selve Sorgen kan være noget Tiltrækkende, sandsynligvis
fordi det stærke Liv, den kraftige Bevægelse,
den Anspænding af Evnerne, som følger med den, i og
for sig medfører en Tilfredsstillelse. Al Kraftanspændelse,
som ikke er uforholdsmæssig, er forbunden med Lyst
følelse. Midt i Sorgen kan da en Lystfølelse ved selve
Sorgens Styrke og Voldsomhed gøre sig gældende, hvilket
viser sig i, at man ikke \emph{vil} udrives af den. Den organiske
Udladning, som er forbunden med enhver stærk
Stemning, er et vigtigt Element her, selv om den ikke
skulde indeholde den fulde Forklaring. --- Da Auguste
Comte havde mistet den Kvinde, der udøvede en saa afgørende
Indflydelse paa hans Aandsretning i den sidste
Periode af hans Liv, udtalte han en Gang i en Henvendelse,
han rettede til sit Erindringsbillede af hende: »Kun
dig skylder jeg, at jeg ikke kommer til at forlade Livet
uden paa værdig Maade at have lært den menneskelige
Naturs bedste Følelse at kende\dots{} Midt under de
tungeste Smerter, Følelsen kan medføre, har jeg aldrig
ophørt at føle, at \emph{det Væsentligste for Lykken} er
at have Hjærtet opfyldt paa værdig Maade, selv om det
er af Smerte, ja selv af den bitreste Smerte«. --- Auguste
Comte og Clärchen ere altsaa ganske enige, og Velfærdsmoralen
er enig med dem. Med den store Livsbevægelse
og Livsfylde følge ogsaa de store Omkostninger, de store
% --- p. 31 ---
\setcounter{footnote}{0}%
\opage{31}Kontraster og de store Svingninger. Det gælder da at
have Mod til at vove og vælge. Hvem har sagt, at det
Højeste skulde kunne faas for en billig Penge?\footnote[1]{Denne sidste Udvikling har jeg allerede fremsat i „Tilskueren“. Nov.---Dec. 1889. p. 861 f.}

Vi have nu én Gang intet andet psykologisk Kriterium
paa Livets Sundhed og Kraft end Lystfølelse. Hvad
der --- i \emph{alle} sine nærmere eller fjernere Virkninger ---
kun volder Lystfølelse hos alle de Væsener, det berører,
det kan umuligt være at forkaste. For saa vidt er Velfærd,
defineret som varig Lystfølelse, det sidste Vurderingsprincip.
Herimod strider det ikke, at det Højeste, vi 
faktisk kende, ofte ved nogle af sine Betingelser medføre
en stor Smerte. Hvo som vil være sikker for Savn og
Sorg, maa ikke elske Noget som helst; men hvis nu
selve Kærligheden var den højeste Lykke?

Smerte er overalt Tegn paa en begyndende Opløsning
af Livet\footnote[2]{Se om denne Teori: „Psykologi“ VI D, 2---3.}. I simpleste Form viser dette sig ved den
»fysiske« Smerte, der opstaar ved Sønderrivning af det
organiske Væv. Det gælder ogsaa for den »aandelige«
Smerte, der bestaar i Sorg, Tvivl eller Anger\dots{} Den
tyder paa en Disharmoni mellem forskellige Kræfter og
Drifter i Sindet, en Disharmoni, der ved at stige tilsidst
kan sprænge Bevidstheden. Naar Smerten dog ofte er
et nødvendigt Gennemgangsled, er det dels, fordi den
(som allerede vist) ved Kontrast og Stimulering kan sætte
stærkere Kræfter i Bevægelse, der opveje de opløsende
Kræfter, dels ogsaa fordi det gælder Opløsning af Noget
i os, der hindrer en friere og rigere Livsudfoldelse. Barnefødsel
er forbunden med Smerte, fordi det nye Liv kun
kan fremtræde ved Løsrivelse fra det gamle, hvoraf det
hidtil har udgjort en Del. Sandhedserkendelsen vindes
ofte med Smerte, fordi Fordomme og Illusioner briste.
% --- p. 32 ---
\setcounter{footnote}{0}%
\opage{32}I Angerens Smerte opløses et lavere Selv, for at et højere
Selv kan udvikle sig.

4. Hvad der især har affødt Modstanden mod Velfærdsprincipet,
er sikkert det, at man er tilbøjelig til
ved Ordet Lystfølelse først og fremmest at tænke paa
de mest elementære og sanselige Former for Lyst. Velfærdsprincipet
udelukker ikke disse, men indbefatter Mere
end dem. Det tager Hensyn til alle Sider i Menneskets
Natur, idet det hævder, at en varig Lystfølelse ikke kan
frembringes med Sikkerhed, naar en væsentlig Side af
denne Natur lades upaaagtet. Manglen ved de elementære
Lystfølelser er, at de for største Delen kun svare
til øjeblikkelige og begrænsede Forhold. Et Væsen, hvis
Følelse er af rent elementær Art, kan hævde sig, saa
længe de simple Livsbetingelser, efter hvilke det er tillempet,
ikke ændres. Saaledes synes nogle af de laveste
Dyreformer, som Infusionsdyr og Rhizopoder, i uendeligt
lange Tider at have bestaaet næsten som de nu ere.
Tillempelsen efter Forholdene er her saa godt som fuldstændig.
Det Samme kan være Tilfældet med Væsener,
der paa et tidligere Udviklingstrin have besiddet mere
udviklede Former og Organer og derfor sandsynligvis
ogsaa have havt et rigere Følelsesliv. Dyr, som i deres
Ungdom leve frit, men senere føre en snyltende Tilværelse,
gaa helt op i Ernæringslivet og miste alle Nerver og
Muskler, der ikke umiddelbart tjene dette. --- Ogsaa hos
Menneskene findes Exempler nok paa en Tillempelse efter
Forholdene, der hindrer videre Udvikling. Om Ildlænderne
bemærker Darwin i sin Skildring af deres (i vore
Øjne) sørgelige Tilværelse\footnote[1]{Rejse om Jorden. p. 242.}: »Der er ingen Grund til at
at tro, at Ildlænderne aftage i Mængde; derfor maa vi % PRINTED AS IS (tro,: at at: at doubled at the line turn (600 ppi))
antage, at de nyde et saadant Maal af Lykke, hvorledes
% --- p. 33 ---
\setcounter{footnote}{0}%
\opage{33}den nu end er beskaffen, at Livet faar Værdi for dem.
Naturen, som gjorde Vanen til en uimodstaaelig Magt
og lod dens Virkninger blive arvelige, har afpasset Ildlænderen
efter hans elendige Lands Klima og Frembringelser«.
Ogsaa højere staaende Naturfolk, og ikke blot
Naturfolk, frembyde lignende Exempler. Det er Nøden,
der driver Mennesket ind paa Fremskridtets Bane. Hvor
 selv usle Forhold ikke have umuliggjort Tillempelsen,
betinges Følelseslivet kun af en snever Kreds af Forhold.
Der kan da maaske opnaas en fuldstændigere Tilfredshed
end under mangfoldigere og mere indviklede Forhold. Et
lille Kar er lettere at fylde end et stort\footnote[1]{Vas minus majore plenius esse potest, qvamvis liqvoris minus contineat. \emph{Cartesius}: Epistola IV. (De vita beata).}.

Det kunde nu synes at maatte være en væsentlig
Indvending mod Velfærdsprincipet, at det maatte lede til
at anerkende en saadan til Livets første Fornødenheder
og til rent elementære Følelser indskrænket Tilværelse
for at staa lige saa højt som den Tilværelse, der udfyldes
med aandeligt Arbejde og Kulturvirksomhed, --- ja, for
at staa højere end denne sidste, der neppe vil medføre
saa ublandet og sikkert Velbefindende, men være forbunden
med stedse ny Trang, ny Stræben og ny Uro.
--- En Tilværelse som Ildlændernes forekommer vel Alle
kummerlig og ussel, fordi de sikkert ofte lide Mangel og
Nød. Men vi kunne ogsaa tage et andet Exempel.
Alexander von Humboldt traf i Sydamerika en Befolkning,
der levede af Banantræerne, som ere saa frugtbare, at
en Tønde Land, som er bevoxet med dem, kan ernære
50 Individer. Disse Træer behøve ingen egentlig Dyrkning;
kun skal Jorden omkring Rødderne bearbejdes med
Hakke en eller to Gange om Aaret. Disse Mennesker
kunne, i Ordets dobbelte Betydning, »nære« sig ved meget
% --- p. 34 ---
\setcounter{footnote}{0}%
\opage{34}lidt Arbejde. Derfor ere de ogsaa sløve og uden Kultur.
Men de Fornødenheder, de have, faa de tilfredsstillede.

Hvad der for os vilde gøre et saadant Liv uudholdeligt,
\emph{Trangen til Virksomhed, Udvikling og}
\emph{Fremskridt}, existerer ikke paa saadanne Trin. Denne
Trang er, hvad der vel maa mærkes, \emph{selv en Følge
af Udvikling og Fremskridt}. --- En af Darwins
Forgængere, \emph{Lamarck}, antog en indre, iboende Udviklingstrang
hos alle levende Væsener. Men \emph{Darwin} har,
støttet til Erfaringen, hævdet, at Udviklingen stedse indledes
ved ydre Forholds Indvirkning. For Lamarck blev
det en Vanskelighed, hvorledes de aller simpleste Former
ere vedblevne at bestaa, --- hvorfor de i Kraft af den
antagne Udviklingstrang ikke forlængst ere gaaede over
til højere Former. Han løste denne Gaade ved at antage
en generatio æqvivoca, en Nyopstaaen af de simpleste
organiske Væsener ud af uorganisk Materie, en Antagelse,
som hidtil ikke er bleven bekræftet. For Darwin volder
de simpleste Formers fortsatte Bestaaen ingen saadan
Vanskelighed. Ti efter hans Lære begunstiges kun de
Variationer, som kunne være til Gavn under Kampen for
Tilværelsen under de givne Forhold; Kvalitetsvalget medfører
derfor ikke nødvendigvis en fremadskridende Udvikling\footnote[1]{Arternes Oprindelse. p. 153 f.}.
For Menneskenes Vedkommende har \emph{Th. Waitz}
klart vist, hvorledes Udviklingstrangen selv er et Udviklingsprodukt:
»Af et Folk uden Samkvem og uden
Væddestrid med andre, et Folk, som med forholdsvis Lethed
eller kun med Overvindelse af tilvante og som nødvendige
betragtede Vanskeligheder kan tilfredsstille sin
Fornødenhed ved hvad der umiddelbart omgiver det, og
som føler sig tilfreds dermed og lever lykkeligt, --- af et
saadant Folk maa man ikke vente, at det skal gøre nogen
% --- p. 35 ---
\setcounter{footnote}{0}%
\opage{35}som helst Anstrængelse for at civilisere sig. Hvo som
har, hvad han behøver, og føler sig vel derved i enhver
Henseende, han arbeider ikke. Man bliver ikke civiliseret
uden Nød og paa fri Haand eller for maaske at følge en
ædel Trang, som bor i Hjertet. Gaar det da anderledes
i vort moderne Samfund? Hører der ikke først en lang
Dannelsens Skole til, for at indgyde Mennesket Interesse
ved Arbeidet som saadant? Man regne en Gang efter
blandt de saakaldte Dannede, hvor mange der gøre sig
Anstrengelse for deres egen eller Andres Dannelses Skyld
uden at være nødte dertil! Man vil saa høre op med at
betvivle de Vildes Jevnbyrdighed med os, hvad aandelige
Evner angaar, i det Mindste af den Grund, at de undlade
at anstrænge sig uden Nødvendighed«\footnote[1]{Die Indianer Nordamerikas. Leipzig 1865. p. 69.}.

Ejendommeligt for Naturtilstanden i Modsætning til
Kulturtilstanden (for saa vidt man skarpt kan stille disse
to Tilstande i Modsætning til hinanden) er det, at hist
maa Stødet til Forandring i Levevis og Tænkemaade
stedse komme ude fra, medens der her --- selv om det
skulde være i forholdsvis smaa Kredse --- rører sig en
stærkere eller svagere Trang til Fremskriden, som stedse
holdes vedlige. Denne Forskel danner en Analogi til
Forskellen mellem uorganisk og organisk Tilværelse. I
Organismen vedligeholdes Kræfternes Spil med en vis
Uafhængighed af Øjeblikkets og de begrænsede Omgivelsers
Paavirkninger; i den uorganiske Verden (og ligeledes
naar man sammenligner de lavere Organismer med
de højere) spille ydre Paavirkninger en langt mere afgørende
Rolle. Hos Kulturfolkene er der vakt en Trang
til at ændre og differentiere, til at udforme Livet i alle
Retninger og til at tilspidse det inden for hver enkelt
Retning. Der er udvoxet aandelige Følehorn, som uaf%
% --- p. 36 ---
\setcounter{footnote}{0}%
\opage{36}brudt ere i Virksomhed. En ny Art Følelser bliver nu
mulig, Følelser, hvis Natur ikke blot bestemmes ved de
Formaal, som opnaas, men ogsaa knytte sig umiddelbart
til selve Arbejdet og Virksomheden. Mennesket er derved
blevet selvstændigere og friere, og hans Aandsliv har
vundet i Dybde og Styrke, idet det ikke længere blot
bestemmes ved den ydre Verden, men væsentligt ved de
Kræfter, der røre sig i den indre Verden. Der stilles
da nu ogsaa højere Fordringer til Livet.

\emph{Denne Udviklingstrangs Opstaaen kan stemme
fuldstændigt med Velfærdsprincipet}. Ti hin Stabilitet
i Naturtilstanden, der snart kan forekomme os
ynkelig, snart paradisisk, er afhængig af de ydre Forholds
Stabilitet. Men en absolut sikret Stabilitet gives ikke i
Naturen. Selv om de nærmeste Omgivelser ikke ændres,
sker der Ændringer andre Steder i Verden og hos andre
Væsener, og Kampen for Tilværelsen kan da nøde de
stabile Væsener til enten at gaa under eller udvikle
Fremskridtsevnen hos sig. De højere udviklede Væsener
stille ofte under Kampen for Tilværelsen de lavere udviklede
Væsener dette Valg. Det stilles mange Naturfolk, eller
Kulturfolk, der ere blevne staaende paa et lavere Trin;
de gaa saa ofte Undergangen i Møde, naar den højere
Kultur nærmer sig.

Hvad her er paavist ved Exempler fra Racer og
Folkeslag, kan ogsaa paavises for de enkelte Individers
Vedkommende. En fuldkommen Tillempelse efter begrænsede
Forhold indeslutter altid en Fare: Faren for, at
Individet, naar dets Livsforhold og dets Synskreds udvides,
ikke skal have de indre Betingelser for at hævde
sig. Den barnlige Naivitet, den drømmende Fantasi,
den sanselige Nydelse have hver sin Ret og sin Betydning;
men Tilstanden bliver let en Søvngængertilstand, i
hvilken der kun er Sikkerhed og Lykke, indtil Opvaag%
% --- p. 37 ---
\setcounter{footnote}{0}%
\opage{37}ningen sker, og hvor den største Hjælpeløshed kan indtræde
ved Opvaagningen. Og der ses endda ganske bort
fra, at saadanne Tilstande ofte kun bestaa paa andre
Væseners Bekostning, idet disse enten maa yde Søvngængeren
Hjælp eller dog undvære hans slumrende Kræfter.

Velfærd kan derfor ikke bestaa i en passiv, én Gang 
for alle frembragt Tilstand, en Tilstand, der ikke selv er
Udgangspunkt for ny Udvikling og Fremskridt. Velfærd
maa i højeste Forstand bestaa i \emph{en Tilstand}, hvor der er
opsamlet Kraft og anlagt rige Muligheder for Fremtiden,
og som afføder en Trang til at sætte nye Formaal og
indlede ny Stræben. Den maa bestaa i en Tilstand,
\emph{der baade i og for sig er attraaværdig og afføder
Stræben efter ny attraaværdige Tilstande,
altsaa en Tilstand, der ikke blot er Formaal,
men ogsaa Middel, ikke blot har Værdi som
Virkning, men ogsaa som Aarsag}. Lystfølelsen
knytter sig her umiddelbart til selve Virksomheden og
Arbejdet, til selve Udfoldelsen af Kræfterne, ikke blot til
Resultatet. Hvor en saadan Lystfølelse er mulig, dér
blive mange Lidelser udholdelige, som paa et lavere
Niveau vilde være Tegn paa selve Livets Opløsning.
Spænding og Længsel, Skuffelse og Savn kunne ikke
udeblive og ville med en vis Rytme afløse den frejdige
Fremgang mod det Maal, man har sat sig; men gennem
alle Svingninger kan Grundretningen og Grunddriften
bevares.

Dette er det Ideal, Velfærdsprincipet, naar det gennemtænkes,
fremkalder for os. Hvor vidt det kan naas, er
et Spørgsmaal, som til hver given Tid og for hvert enkelt
Individ maa besvares ad experimental Vej. Det gælder
meget ofte ikke blot om Ændringer i Individernes Natur,
men ogsaa i Samfundets Ordning eller endog i den ydre
Natur, for at en Tilnærmelse til Maalet kan finde Sted.
% --- p. 38 ---
\setcounter{footnote}{0}%
\opage{38}Men for os er i denne Sammenhæng Hovedsagen, at vi
her have en Maalestok til Bedømmelse af Handlinger og
Institutioner, % PRINTED AS IS (Institutioner,: a comma ends the paragraph (600 ppi))

Denne Maalestok svarer til en Tendents, der gaar
gennem hele den organiske Natur, hvor Lyst i det Hele
og Store tyder paa Liv og Fremgang, Smerte paa Tilbagegang
og Død. Velfærdsprincipet hævder \emph{Livets Ret}:
ethvert Væsen har Ret til at bestaa og udvikle sig, for
saa vidt derved ikke frembringes saa meget des større
Smerte for andre Væsener. Den Etik, der bygger paa
Velfærdsprincipet, søger altsaa at fortsætte Naturudviklingen
paa mere bevidst og harmonisk Maade. Den opfordrer
til at søge efter Midler til Liv og Fremgang, som
den ubevidste Naturudvikling ikke vilde føre til, og den
stræber at udelukke den unødvendige Smerte og Nød,
som den blinde Kamp for Tilværelsen fører med sig. Den
omfatter en Række Opgaver, der strækker sig lige fra
den Lindring og Hjælp, som Medlidenheden motiverer, til
Arbejdet paa de højeste sociale, æstetiske og intellektuelle
Idealers Virkeliggørelse. Det er den specielle Etiks Opgave
at undersøge, hvorledes disse forskellige Opgaver
skulle behandles, og hvorledes de forholde sig til hverandre
indbyrdes.

5. Men fordi Velfærd bestaar i Virksomhed og Udvikling,
kan man ikke omvendt sige, at Virksomhed og
Udvikling nødvendigvis ere forbundne med Velfærd eller
føre til Velfærd. Fordi det er urigtigt at tro, at Begrænsning
af Fornødenhederne altid fører til Maalet, følger jo
dog ingenlunde, at ubegrænset Mangfoldiggørelse af Fornødenhederne
skulde være det Rette. Kulturudviklingen
kan antage Former og slaa ind paa Veje, der ikke stemme
med Velfærdsprincipet. Vi finde derfor ogsaa i Historien
til sine Tider skarpe og indtrængende Advarsler mod den
bestaaende Kultur. Saaledes i Grækenland af Sokrates,
% --- p. 39 ---
\setcounter{footnote}{0}%
\opage{39}Kynikerne og Stoikerne, i det 18de Aarhundrede af Rousseau
og i vore Dage af Leo Tolstoj. Saadanne store
Aanders Opposition maa i hvert Tilfælde gøre os opmærksomme.


Jeg ser her bort fra det Spørgsmaal, om det, vi kalde
Kultur, uden videre kan bibringes ethvert Folk. Selv
om man har været for hurtig til at frakende Naturfolk
Evnen til at modtage Civilisation, følger deraf ikke, at
det vilde være et Gode for et Folk at opgive sine i og
ved det selv udviklede Livsformer for at lade sig regulere
efter Former og Idealer, der ere udviklede under helt
andre Forhold. Saa direkte kan selv den mest berettigede
og fuldkomne Kultur ikke meddeles. \emph{Waitz}, der udtrykkeligt
hævder, at man intet Bevis har ført for Indianernes
Uimodtagelighed for Kultur, priser dog de Indianerhøvdinger,
der modsatte sig den europæiske Civilisations
Indtrængen i deres Folk, for »Kærlighed til deres Folk
og rigtig Indsigt i dets sande Velfærd«. Der ligger her
(lige som ved etiske Systemers Sammenstød) en Mulighed
for tragiske Konflikter, som særligt ogsaa Indianernes
Historie frembyder mange Exempler paa.

Grunden til, at der kan komme Strid mellem Kultur
og Velfærd, ligger i Kulturbestræbelsens Rastløshed og
Hensynsløshed. Det gaar med den som med den uvilkaarlige
(»spontane«) Bevægelsestrang, der fører til at
bruge Kræfter og Lemmer, blot fordi der Energi til Stede, % PRINTED AS IS (Energi: fordi der Energi: er missing (600 ppi))
uden at Henblikket til et værdifuldt Formaal leder den,
saa at det bliver tilfældigt, hvilket Resultat Bevægelsen
faar. Kulturbestræbelsen kan føre dels til Overanspændelse
af Kræfterne, dels til ensidig Anvendelse af dem
(formedelst overdreven Differentiation), dels til Isolation,
brudstykkeagtig Udskillen af de enkelte Virksomheder.
Hos det enkelte Individ plejes og udvikles enkelte Evner
(hos Nogle f. Ex. Intelligentsen, hos Andre den fysiske
% --- p. 40 ---
\setcounter{footnote}{0}%
\opage{40}Arbejdskraft) paa andres Bekostning; Harmonien, Evnen
til at føle sig som Helhed og Enhed, mangler. En saadan
Ensidighed gør, at Individet kun faar Værdi som et
Hjul i det store Maskineri, kun bliver Middel, ikke Formaal.
Den hænger sammen med en ensidig social Udvikling;
Undertrykkelsen af enkelte Sider i den Enkeltes
Natur, foraarsages ved Undertrykkelsen af enkelte Stænder
og Klasser i Samfundet. Dersom man vilde identificere
Kultur og Etik og betragte Udvikling som et sikrere
Kriterium end Velfærd, vilde man ganske overse, at der
gives et socialt Spørgsmaal. Det sociale Spørgsmaal er
et etisk Spørgsmaal og tillige et Spørgsmaal om vor
Kulturs Berettigelse, begge Dele i Kraft af Velfærdsprincipet.
Skulde det være det Rette, at den materielle
og den ideelle Kulturs Frembringelser kun komme et lille
Mindretal til Gode, medens alle Andre udelukkes fra at
faa Del i dem? Det vilde ganske stride mod det Ideal
af et Samfund, som følger at Velfærdsprincipet. Ti den % PRINTED AS IS (følger: følger at: at for af (600 ppi))
største Velfærd er til Stede, hvor hvert enkelt Individ
udvikler sig paa selvstændig Vis saaledes, at det derved
paa bedste Maade hjælper andre Individer til en lignende
Udvikling ud fra deres ejendommelige Stade. Der existerer
da \emph{et harmonisk Samfund af selvstændige Personligheder}.
Dette er nu den højeste etiske Ide. Hvert
enkelt Individ staar som en lille Verden for sig, og dog
i den inderligste Vexelvirkning med den store Verden,
af hvilken det er en Del. Den Enkelte tjener Slægten,
og Slægten tjener den Enkelte. Enhver Særstilling, enhver
Ulighed i Fordelingen maa begrundes ved hvad de
forskellige Livsvilkaar og Opgaver kræve, og ethvert Individs
Evner og Drifter skulle udvikles saa fuldt og rigt, som
foreneligt er med hvad Slægtens Liv som Helhed kræver\footnote[1]{Se nærmere herom „Etik“. p. 31. 164. 227---240.}.

% --- p. 41 ---
\setcounter{footnote}{0}%
\opage{41}6. Det følger af hvad her er udviklet, at Velfærdsprincipet
ingenlunde er saa let at anvende i de enkelte
Tilfælde, som det maaske kunde synes. Den specielle
Sammenhæng kan være saa indviklet, at vi ikke kunne
overse den eller forudse Virkningerne af vor Indgriben.
Vi kunne ikke af Velfærdsprincipet a priori udlede et
System af enkelte Handlinger eller en Ordning af Samfundet
eller en hel Kultur. Dets Betydning er (lige som
Aarsagsprincipet paa det teoretiske Omraade) at stille
Problemer og at vejlede ved deres Behandling. Det er
regulerende, ikke konstruerende. Det forudsætter det
umiddelbare, uvilkaarlige Liv hos den Enkelte og i Samfundet
og træder først til, naar der sker en bevidst
Drøftelse af det saaledes Udvikledes Værdi og af Maaden,
hvorpaa det skal føres videre. Al Etik faar derved en
historisk Karakter. Ti vi begynde aldrig --- hverken hvad
den enkelte Individualitet eller Samfundet angaar --- absolut
forfra, men have at arbejde videre paa et givet Grundlag
under Vejledning af de Principer og Idealer, der udspringe
af vor Natur.

% --- p. 42 ---
\setcounter{footnote}{0}%
\essay[Forholdsloven i Etiken]{Forholdsloven i Etiken.}
\opage{42}


\begin{center}I.\end{center}


Det er Hensigten med denne Afhandling at paavise,
at Gyldigheden af alle etiske Love beror paa bestemte
Forhold og Betingelser, saa at Forholdsloven, vor Erkendelses
Grundlov, ogsaa finder Anvendelse paa Etikens
Omraade. Det vil vise sig, at denne Betragtning i væsentlige
Henseender tjener til at karakterisere den etiske
Erkendelses og de etiske Loves Beskaffenhed. ---

Vi kunne begynde med at undersøge, i hvilken Betydning
og med hvilken Ret, man i det Hele taler om
etiske Love. Ordet Lov bruges her ikke i samme Betydning
som i Naturvidenskaben, i Psykologien og i
Sociologien. Ti Etiken vil ikke paavise de Regler, som
den menneskelige Villen og Handlen faktisk følger,
men de Normer og Idealer, den \emph{skal} følge, den Maalestok,
den bør være underlagt. Den Betydning af Ordet
Lov, som ligger til Grund i den filosofiske Etik, minder
derfor om den Maade, paa hvilket Ordet bruges i den
positive Retsvidenskab og desuden i den teologiske Etik,
saa vel som i Videnskaben om den positive Moralitet
(de i Sæd og Skik og i den bestaaende offentlige Mening
liggende Normer). Men fra disse Videnskaber afviger
den filosofiske Etik ved, at den ikke bygger paa nogen
% --- p. 43 ---
\setcounter{footnote}{0}%
\opage{43}ydre, naturlig eller overnaturlig Autoritet, men søger at
grunde sine Normer og Idealer paa selve den menneskelige
Natur, saaledes som vi erfaringsmæssigt kende denne.
Man har ment, at denne Omstændighed gjorde det utilstedeligt
at tale om etiske Love.

I Følge \emph{John Austin} forudsætter en Lov to
Parter, en, som giver Loven, og en, for hvem Loven
gælder, og Lovens forbindende Kraft ligger i, at den,
der giver Loven, har Magt til at hævde den ved at
knytte Smerte til dens Overtrædelse. Da nu den saaledes
med Magt hævdede Lov frembringer Ensformighed i
Handlemaaden, idet de enkelte Handlinger bestemmes
ved den i Loven udtalte Fordring, kan Ordet Lov tillige
komme til at betyde en almindelig Regel, som følger i
de enkelte Tilfælde, og det er denne Side af Begrebet
Lov, og kun denne, som vi gøre Brug af, naar vi tale
om Naturlove. Vi mene dermed den regelmæssige Rækkefølge
af Naturens Fænomener, uden at vi behøve at forestille
os denne Regelmæssighed som Virkning af en bydende
Villies Magt. Selv om nu denne Sprogbrug kan
tilstedes, negter Austin bestemt, at Begrebet Lov kan
finde Anvendelse paa de Regler for Handlen, som det
enkelte Menneske fastsætter og anerkender for sig selv.
Ti her er der jo hverken Tale om en Naturlov, ej
heller gives der en over Individet staaende Magt, som
kan hævde Loven i Overtrædelsestilfælde. »Ti skønt han
kan fatte det oprigtige Forsæt at tilføje sig selv Smerte,
hvis hans Færd afviger fra den Rettesnor, han har fastsat
for den, saa afhænger den virkelige Tilføjelse at % PRINTED AS IS (Tilføjelse: Tilføjelse at: at for af (600 ppi))
denne mulige Smerte af hans egen Villie«)\footnote[1]{Lectures ou Jurisprudence. I. p. 213 f.}.

Skulde der dog nu ikke være et \emph{psykologisk
Grundlag for Begrebet etisk Lov}, der kunde inde%
% --- p. 44 ---
\setcounter{footnote}{0}%
\opage{44}holde de tre Elementer, som i Følge Austin forudsættes
ved den strenge Anvendelse af Begrebet Lov: 1) et Forhold
mellem Overordnet og Underordnet, 2) en almindelig
Regel for Handlen og 3) Hævden af denne almindelige
Regel ved Smerte i Overtrædelsestilfælde?

Der rører sig hos Mennesket forskellige Følelser,
hver bestemte ved sit Fornemmelses- eller Forestillingsindhold.
Der er en Følelse, som er knyttet til Individets
egen Selvopholdelse. Og der er andre Følelser, som lader
Individet deltage i Andres Ve og Vel. Da nu enhver
Følelse tillige er et Handlingsmotiv eller afføder en Drift
til Handling, vil der i Menneskets Bevidsthed kunne
fremtræde en Drift til Handling, som er bestemt ved
Hensynet til en mere omfattende Tingenes Orden, til en
større Gruppes Velfærd, over for en Drift til Handling,
som kun er bestemt ved Hensynet til Individets egen
begrænsede Interesse. Der vil da i Individet være Noget,
som peger ud over dette selv og anviser det Plads som
et enkelt Led i et stort Hele, saaledes at dettes Velfærd
under Tiden kan kræve Opofrelse af, hvad Hensynet til
Individets egen Velfærd som isoleret Væsen vilde fordre.
Her vil der da i Individets eget Væsen være en Modsætning
mellem to Parter, mellem en Drift, som er knyttet
til Forestillingen om en større Tingenes Orden, og en
Drift, der har en snevrere Horisont. Det første Element
i Lovens Begreb er hermed paavist. --- I Virkeligheden
forholder det sig ikke anderledes, selv hvor den Part,
der giver Loven, er en ydre Autoritet. Ti denne Autoritet
maa dog kundgøre sig for Individet ved at vække
dets Følelse (hvad enten denne Følelse er Frygt eller
Ærefrygt), og Loven vil ogsaa her lægge sig for Dagen
gennem et Forhold mellem to Følelser eller Drifter i
Menneskets Indre. For den, som kan udrydde Følelsen
af Frygt eller Ærefrygt i sig, existerer Loven ikke. Det
% --- p. 45 ---
\setcounter{footnote}{0}%
\opage{45}ses heraf, at den filosofiske Etik har lige saa megen Ret
til at anvende Begrebet Lov som Retslæren, Teologien
og den positive Moral. --- Jeg gaar her ikke nærmere
ind paa, hvorledes denne indre Modsætning mellem et
mere omfattende og et mere begrænset Hensyn inden for
Følelses- og Driftlivet kommer i Stand. Det er en Modsætning,
der historisk har udviklet sig under det sociale
Livs Indflydelse paa Mennesket i Kraft af almindelige
psykologiske Love.

Altsaa det første af de tre Elementer i Lovens Begreb
kan meget vel genfindes i de Love, den filosofiske
Etik opstiller. Det Samme gælder nu ogsaa om det af
de tre Elementer, som Austin aller mindst kunde finde
her, nemlig Sanktionen eller den ved Overtrædelse af
Loven hidførte Lidelse. Austin mener, at det blot beror
paa Individets Villie, om det kommer til at lide Smerte
ved Overtrædelsen af den Lov, der skulde udspringe af
dets egen Natur. Men naar den etiske Lov er Formuleringen
af en Drift, som fremkaldes ved, at Individets
Følelse er bestemt ved Hensynet til et større Hele, da
vil en saadan Drift ikke kunne møde Modstand og Modsigelse
fra andre Sider i Individets Natur, uden at der
opstaar mere eller mindre stærk Smerte. Naar den Drift,
der volder, at Individet ikke lyder den etiske Lov, det
selv har anerkendt, ikke helt kan fortrænge eller udrydde
den Drift, som førte til Anerkendelse af Loven, saa vil
der opstaa den indre Disharmoni, den smertelige Følelse
af indre Opløsning og Nedværdigelse, som vi kalde ond
Samvittighed eller ogsaa Anger. Og det beror ikke paa
Individets vilkaarlige Valg, om denne Smerte skal indtræde
eller ikke; den følger med psykologisk Nødvendighed
af de forhaandenværende Betingelser. Det, der ytrer
sig i den onde Samvittighed og i Angeren, er ikke noget
Tilfældigt eller Vilkaarligt. Det er et under de givne
% --- p. 46 ---
\setcounter{footnote}{0}%
\opage{46}Forhold uundgaaeligt psykologisk Resultat --- lige saa
uundgaaeligt, som den ydre Straf er, naar Statsmagten
er stærk, klog og hurtig nok til at faa fat i Overtræderen
af den borgerlige Lov. Men Statsmagten formaar
jo ikke altid at ramme Forbryderen med sin Straf, og
saaledes kan det ogsaa ske, at Angeren udebliver, nemlig,
naar den Drift, der fører til Anerkendelsen af et
større Hele, er udslukt, eller naar den ikke er stærk og
levende nok til at gøre sig gældende i Modsætning til
den Drift, som førte til at udøve den Handling, der stred
imod hin Anerkendelse.

Endeligt vil ogsaa den Egenskab ved Loven, at den
betegner en almindelig Regel, kunne paavises i den etiske
Lov. I Forestillingen om den større Helhed, til hvilken
Mennesket føler sig knyttet, --- Familien, Klanen, Nationen,
Menneskeheden, --- har det et Princip, af hvilket
der kan udledes almindelige Regler for dets Villen og
Handlen. Naar Individet ikke vil sige sig selv imod,
maa det bedømme sin Villies Tilstand og Retning ud fra
det mest universelle og ophøjede Synspunkt, det kender,
ud fra Tanken om det store Hele, af hvilket det føler
sig som et enkelt Led. De enkelte Regler følge med logisk
Konsekvents ud fra denne Tanke, naar denne én
Gang er anerkendt. Den etiske Videnskab er da kun en
Udfoldelse af, hvad der implicite er givet i den universelle
Sympati, saaledes som denne historisk har udviklet
sig i Menneskeslægten. ---

Men naar det saaledes har vist sig berettiget at anvende
Begrebet Lov i Etiken, saa ligger der tillige i den
foregaaende Begrundelse, at saa vel det almindelige
Princip som de Regler, der udledes af det, kun gælde
i \emph{Forhold til} og \emph{paa Betingelse af} et bestemt psykologisk
Grundlag. Dersom der skulde gives Mennesker,
som mangle dette Grundlag, saa kunde og vilde de ikke
% --- p. 47 ---
\setcounter{footnote}{0}%
\opage{47}anerkende de etiske Love. Hvorledes man maa stille sig
overfor saadanne Individer, om de findes, er en Sag for
sig, som vi senere ville komme til at berøre; her hævde
vi kun, at alle etiske Love kun gælde i Forhold til en
bestemt psykologisk Forudsætning. Dette er Etikens \emph{psykologiske}
Relativitet\footnote[1]{Smlgn. min „Etik“. p. 35---37. --- Se ovenfor p. 7---17.}.

\begin{center}II.\end{center}


Men det er ikke tilstrækkeligt at fastsætte det psykologiske
Grundlag og opstille det dertil svarende etiske
Princip. De Forhold, under hvilke de af dette Princip
følgende Regler skulle anvendes, ere saa indviklede og
sammensatte, at det er vanskeligt, ja umuligt at finde
Regler, som uden videre skulde kunne finde Anvendelse
i alle \emph{specielle Forhold og Situationer}.

Den \emph{almindelige} Lov eller Regel kan ikke forudsé
alle mulige Situationer, i hvilke et Menneske kan blive
nødt til at fatte en afgørende Beslutning. Det er en
Vanskelighed, som den etiske Lov har fælles med den
juridiske og den teologiske Lov, saa vel som med den
positive Moralitets Lov. Livet er for rigt og mangfoldigt
til at kunne sættes i et fuldt færdigt System. Forholdene
gaa gennem mange forskellige Nuancer og Mellemtrin
over i hverandre; skarpt afstukne Grænser existere ikke.
Den almindelige Lov, der skal passe paa \emph{alle} Tilfælde
af en vis Art, bliver let uanvendelig i det enkelte
bestemte Tilfælde, som kan være aldeles enestaaende.
Den etiske Lov maa derfor --- hvis den ikke for den tilsyneladende
formelle Konsekvents vil ofre Livets virkelige
Ret --- fastsætte den almindelige \emph{Tendents} eller \emph{Retning},
som Villien skal følge, og ikke give specielle An%
% --- p. 48 ---
\setcounter{footnote}{0}%
\opage{48}visninger  til, hvorledes der i alle enkelte Tilfælde ubetinget skal handles. Kun paa denne Maade kan Lovens
Almengyldighed forenes med de til forskellige Tider og
paa forskellige Steder forskellige Forhold. Dette kan man
kalde \emph{den historiske Relativitet} af de etiske Loves
Indhold\footnote[1]{Se min „Etik“. p. 160 f.}. Ud fra samme psykologiske Grundlag og i
Henhold til samme almindelige Vurderingsprincip vil der
under forskellige historiske Forhold, i forskellige Situationer
være at træffe forskellige Afgørelser med Hensyn
til Samfundets Ordning og det enkelte Individs egen
Færd. Hvad der er muligt paa et Udviklingstrin, er ikke
muligt paa et andet. Der kan ikke af det almindelige
Velfærdsprincip (saa lidt som af noget andet Princip)
konstrueres nogen Ordning af Staten eller Familien, der
uden videre kunde indføres overalt. Enhver social Ordning
forudsætter, at de, der skulle være Samfundets Medlemmer,
leve under visse bestemte indre og ydre Betingelser.
Paa bar Grund kan ingen Bygning opføres.

De etiske Love miste herved ikke deres Betydning.
De kunne ikke sige, hvorledes vi paa Livets Hav skulle
kunne klare hver enkelt Bølge, der rejser sig overfor os,
eller hvorledes vi skulle krydse i de enkelte Tilfælde.
Men de kunne sige os, i hvilken Retning vi i det Hele
skulle styre, og de gøre det muligt, at vort Kompas kan
lede os uden Misvisning, selv om Vejen er nok saa
bugtet og uregelmæssig. De formulere de almindelige
Fordringer, som det højeste Formaal, vi anerkende, stiller
til os. Midlerne til at tilfredsstille disse Fordringer kunne
skifte indenfor vidde Grænser. Men naar Villiens almindelige % PRINTED AS IS (vidde: for vide (600 ppi))
Retning er given, og naar der er tilstrækkelig
Kraft og Konsekvents i Karakteren, ville ogsaa de enkelte
Veje og Midler udfindes, uden at man kommer i
% --- p. 49 ---
\setcounter{footnote}{0}%
\opage{49}Konflikt med Hovedformaalet. \emph{Kant} har med Rette sagt,
at Kasuistik ikke er en Videnskab, men en Kunst, »ikke
en Lære om, hvorledes Noget skal findes, men Øvelse i
hvorledes Sandheden skal søges«.

Kasuistiken angaar de Problemer, som de givne objektive
Forholds indviklede Natur stiller, naar der skal
tages en med de etiske Principer stemmende Afgørelse.
Den søger at analysere Forholdene i deres Elementer for
at udfinde, hvad Konsekventsen af de almindelige Principer
maa være under disse bestemte Forhold. Den opstiller
almindelige Regler for Konflikttilfælde eller søger
ved udtænkte Exempler at øve Evnen til at finde det
Rette. Men Analyse og almindelig Regel kunne ikke
altid overføres fra det ene Tilfælde til det andet. Det
kan f. Ex. være rigtigt, at et højere, mere omfattende
Formaal skal fyldestgøres fremfor et lavere, mere begrænset;
men Vanskeligheden er netop at afgøre, hvad der
i \emph{dette bestemte} Tilfælde er det mest betydningsfulde
Formaal. Afset fra Tilfældenes varierende Mangfoldighed
hænge alle Formaal saaledes sammen, at man ikke kan
gribe ind i et af dem, uden at det faar Virkninger for
en stor Mængde af de andre; og hvorledes da med
Sikkerhed gøre Regnestykket op om, hvad der vil vise
sig at have Overvægt i det Hele, de gavnlige eller de
fordærvelige Virkninger af en vis Art Indgriben?

 Man kan ikke uden videre sige, at det mere omfattende
Hensyn i det enkelte Tilfælde skal have Overvægten
over det begrænsede Hensyn, saa at f. Ex. hvor
Statens og Familiens Interesser komme i Strid, Familien
altid bør vige, eller min Interesse altid skal underordnes,
hvor den kommer i Strid med flere Andres. Ti Livets
sunde Udvikling beror ikke blot paa dets Omfang, men
ogsaa paa dets Kraft og Inderlighed, ikke blot paa Extensiteten,
men ogsaa paa Intensiteten. I det snevrere
% --- p. 50 ---
\setcounter{footnote}{0}%
\opage{50}Samfund kan der leves et inderligere og fyldigere Liv
end i det store Samfund, og selve den Kraft, der holder
det store Samfund sammen, udfoldes og øves i de mindre
Samfund, saa at hints Livskraft beror paa disses. Indenfor
det snevrere Samfund opstaar en Samfølelse, der efter
at have voxet sig stærk kan udvides til at blive en
væsentlig Del af det, der bærer de større Samfund. Den
egentlige Produktivitet finder Sted i de mindre Kredse;
fra dem udgaa de nye Tanker og de nye Kræfter. Derfor
kan Hensynet til Omfanget ikke uden videre blive
afgørende i Konflikttilfælde. Noget Lignende gælder,
hvad det enkelte Individs Selvopholdelse angaar. Slægtens
Liv vinder i Kraft, jo mere og jo bedre de enkelte Individer
formaa at opholde og hævde sig selv saa vel i
materiel som i aandelig Henseende. Inden for den individuelle
Etik skal Selvhævdelsen altsaa ikke uden videre
ofres for Hengivelsen, lige saa lidt som inden for den sociale
Etik det snevrere Samfunds Interesse altid skal
ofres for det størres. --- Men hvor vanskeligt kan det da
ikke i det enkelte Tilfælde være at træffe den rette Afgørelse,
saa at hverken Livets Kraft og Inderlighed eller
dets Omfang og Udbredelse lider Skade? ---

Et andet Tilfælde, i hvilket der kan være Brug for
Kasuistik, er Forholdet til andre Menneskers Fordomme.
Her har Bentham\footnote[1]{Of the influence of time and place in matter of legislation. (Works. I.). p. 180.} givet følgende Regel: »Den, som
letfærdigt og uden Nødvendighed angriber en Fordom, og
den, som finder sig i at blive ledet som en blind Slave
af Fordom, fare begge vild fra Fornuftens rette Vej«. ---
Men hvorledes kan man i det enkelte Tilfælde afgøre,
om der er Nødvendighed til Stede for et Angreb, eller
om det vilde være uværdigt Slaveri at lempe sig efter
% --- p. 51 ---
\setcounter{footnote}{0}%
\opage{51}de herskende Fordomme? Her variere Forholdene i en
særligt høj Grad, da Spørgsmaalet er, om de enkelte
Individer, vi staa overfor, have Betingelser for at kæmpe
den (maaske baade ydre og indre) Kamp igennem, i
hvilken vi styrte dem ind ved at angribe, hvad der hidtil
har staaet urokkeligt fast for dem, og hvad der maaske
har givet deres Liv Hold og Styrke. ---

Bentham mener vel, at man i vore Dage er for tilbøjelig
til at benegte, at de Love, som for vor Tid ere
de bedst mulige, ogsaa vilde have været det i tidligere
Tider. Der gives Love, som vilde have været gavnlige
til enhver Tid. Men han udtaler den historiske Relativitets
Princip paa en træffende Maade i følgende Ord\footnote[1]{Anf. Skr. p. 193.}:
»Hvis jeg skulde vælge, hvad jeg med størst Sandhed
og Redebonhed vilde tillægge de prægtige Egenskaber
Almindelighed og Uforanderlighed, saa skulde det snarere
være visse \emph{Grunde} for Lovene end Lovene selv: snarere
de \emph{Grundsætninger}, paa hvilke Lovene skulle bygges,
de underordnede \emph{Betragtninger}, som følge af disse
Grundsætninger, den bedste Plan, efter hvilken de kunne
forbindes, de \emph{Hensyn}, ved hvilke Lovgiveren bør lade
sig lede, --- end hvilke som helst bestemte Love, som
Lovgiveren burde give, for at de, der ere betroede til
hans Omhu, kunde blive regerede vel«.

\begin{center}III.\end{center}


1. Foruden \emph{det psykologiske Grundlag} og
\emph{den historisk givne Situation}, er der endnu en
tredie Faktor at tage Hensyn til, et tredie Forhold, ved
hvilket den etiske Afgørelse i de enkelte Tilfælde maa
bestemmes. Det er dette tredie Punkt, som det særligt
% --- p. 52 ---
\setcounter{footnote}{0}%
\opage{52}var min Agt ved denne Afhandling at henlede Opmærksomheden
paa, fordi det er blevet mindre paaagtet end
de to første (især end det andet). Det fremkommer ved,
at forskellige Individer paa samme Grundlag (og derfor
med samme etiske Principer) og i de samme Situationer
møde med \emph{forskellige Anlæg og Evner}. Lige saa
lidt som samme Principer føre til samme Konsekventser
i forskellige Situationer, lige saa lidt kunne de i samme
Situation føre til samme Konsekventser, naar de villende
og handlende Individer ere forskelligt udstyrede. Vi staa
her overfor \emph{den individuelle Relativitet} eller den
personlige Ligning i Etiken\footnote[1]{I min „Etik“ p. 133---134 har jeg allerede fremhævet dette Punkt. Bestemtere er det udtalt i den tyske Udgave (Leipzig 1888). p. 154 f., hvor jeg har betegnet Loven om den individuelle Relativitet som Forholdsloven i Etiken. Begrebet Forholdslov kan dog gives en videre Betydning, saa at ogsaa den psykologiske og den historiske Relativitet gaa ind derunder.}.

De forskellige etiske Opfattelser synes i det Hele at
være enige om, at de etiske Love eller Bud ere de
samme for Alle i samme Situationer. Den Etik, som lægger
Autoritetsprincipet til Grund, gaar ud fra, at Autoritetens
Villie forlanger det Samme af alle Individer. Paa den
ene Side staar den evige Villie, paa den anden Side
Menneskene, der skulle lyde dens Bud. Paa lignende
Maade ere Statens Love og de Bud, Sæd og Skik indskærpe,
de samme for alle Individer. Der gælder, som
man siger, ingen Personsanseelse. I Følge Fornuftmoralen
udspringer der en almindelig Lov af den menneskelige
Fornuft, som Enhver maa følge med Fare for at komme
i Strid med sig selv, det vil sige: med det Almenmenneskelige
i sin Natur. Ogsaa Velfærdsprincipets Tilhængere
have i Regelen været tilbøjelige til at opfatte Sagen paa
lignende Maade. De etiske Love udlede de ved Deduk%
% --- p. 53 ---
\setcounter{footnote}{0}%
\opage{53}tion fra Velfærdsprincipet, sammenholdt med den givne
Situation, og disse Love skulle saa lyde ens for Alle.

At der maa være en Fejl i denne Antagelse, kan
indses paa følgende Maade. \emph{Paa Grund af Individernes
forskellige Evner og Drifter, vil en og
samme Fordring faa højst forskellig Betydning}
for dem. Fordringen om en vis Grad af Selvbeherskelse
overfor en eller anden Naturdrift vil saaledes kunne betyde
noget højst Forskelligt for de forskellige Individer,
idet enhver af disse kan have Naturdriften anlagt i forskellig
Grad. Hvad den Ene formaar uden mindste Anstrengelse,
fordi der er en umiddelbar Harmoni mellem
Fordringen og hans oprindelige Natur, formaar den Anden
maaske kun efter en alvorlig, sønderslidende Kamp, medens
den Tredie maaske aldeles ikke formaar at overvinde
sin ubændige Natur. Hvad der for den Første
stemmer med selve Naturtilbøjeligheden, bliver for den
Anden en med stærk Anspændelse og bevidst Villen udført
Pligthandling, og betegner for den Tredie en Opgave,
han ikke kan løse, da han fra først af mangler
Betingelserne derfor. Loven er altsaa i \emph{Virkeligheden}
ikke den samme for dem alle tre. Og nu kan det ikke
nytte at skelne mellem Loven i sig selv og Individernes
forskellige Evne til at opfylde den. Den etiske Lov adskiller
sig jo fra en teologisk eller juridisk Lov ved at
udspringe af selve den menneskelige Natur. Den kommer
ikke udefra, men skal udfolde sig for Individet ud af det
Inderste i dets Væsen. Den kan da ikke naas blot ved
Deduktion ud fra Principer og Situationer, som ere fælles,
og som ikke indeholde noget Hensyn til det enkelte Individs
Ejendommeligheder. Saa længe kun hine to Punkter
ere medtagne ved Lovens Deduktion, kan den ikke være
Individets egen Lov, men en abstrakt upersonlig Fordring.
\emph{Den virkeligt etiske Lov maa ikke blot}
% --- p. 54 ---
\setcounter{footnote}{0}%
\opage{54}\emph{være rettet til Individet, men ogsaa være individualiseret
saaledes, at Individets Væsen kommer
til Udvikling gennem Fyldestgørelse af
Lovens Fordring}. Den maa tage Hensyn til de specielle
Udgangspunkter for Villen og Handlen, som ligge
i det enkelte Individs Natur. Den vilde ellers fordre
Andet, og Mere eller Mindre, af det ene Individ end af
det andet: kun naar den lyder forskelligt for hvert enkelt
Individ alt efter dets Muligheder, stiller den i Virkeligheden
den samme Fordring til dem.

Dette følger simpelt hen af Velfærdsprincipet. Det
enkelte Individ er ikke blot underlagt den Lov, som udtaler,
hvad Slægtens Velfærd fordrer, men er selv et ejendommeligt
Led af Slægten. I den Lov, som gælder for
det, maa derfor to Hensyn mødes paa harmonisk Maade:
Slægtens og Individets Krav. Loven er ikke funden, er
ikke den rette Lov, saa længe ikke begge disse Hensyn
ere komne til deres fulde Ret. Idet den Enkelte i sin
Stræben og Virken arbejder for Slægtens Velfærd, maa
han tillige arbejde for sin egen Velfærd som Del at % PRINTED AS IS (tillige: Del at Slægten: at for af (600 ppi))
Slægten. Derfor maa han have sin specielle Etik, der
baade kvalitativt og kvantitativt er forskellig fra de
andre Individers, idet hans Natur og Dispositioner blive
medbestemmende for den Maade, paa hvilken Loven
former sig for ham.

Det vilde være en stor Misforstaaelse, hvis man
heraf vilde uddrage den Slutning, at al Selvopofrelse var
at forkaste. Den Selvopofrelse, der udspringer af begejstret
Kærlighed, af fri Beslutning, er netop en Ytring af Individets
eget inderste Væsen, Afløb for en dyb Trang
hos det. I den ideale Selvopofrelse gaa Selvhævdelse og
Hengivelse i Et: Livet faar, saa underligt det lyder, kun
Værdi ved at ofres. (Se p. 23. 29ff.). Spørgsmaalet er netop
i de enkelte Tilfælde, om der er Evne hos Individet til
% --- p. 55 ---
\setcounter{footnote}{0}%
\opage{55}saadan Opofrelse. Den Art Handlinger kan lige saa lidt
som andre Arter af Handlinger fordres med lige stor
Lethed af enhver. Her kommer Problemet frem, da Evnen
og Driften til Begejstring og Hengivelse faktisk ikke er
den samme hos Alle.

Lige som i Følge Fechners Lov Fornemmelsen ikke
beror paa Indtrykkets absolute Styrke, men paa Forholdet
mellem det nærværende Indtryk og det forudgaaende,
som allerede har paavirket det sansende Individ,
saaledes maa fra et rent etisk Synspunkt den Fordring,
der i Virkeligheden stilles til Individet, ikke være Et
med den Fordring, der --- uden Hensyn til Individets
Forudsætninger --- kunde udledes af de almindelige Principer
og den givne Situation, men bero paa Forholdet
mellem denne rent objektive Fordring og Individets Forudsætninger
eller Dispositioner, og dette Forhold vil være
forskelligt for de forskellige Individer. Den rent objektive
Fordring kan sammenlignes med en Byrde, der i sig
selv har samme Vægt, men som mærkes forskelligt paa
forskellige Individers Skuldre.

Det er altsaa en fuldstændig Individualisation af den
etiske Lov, som herved fordres. Den Formulering, som
kan passe for Alle, eller dog for et stort Antal Individer,
kan kun være grov og overfladisk. Ved Bestemmelsen af
$\pi$ kunne vi under Tiden lade os nøje med $\frac{22}{7}$, medens vi
ved at stille større Fordringer til Nøjagtighed komme til
flere Hundrede Decimaler, og endda stedse kunne tænke
os større Nøjagtighed. Lige som der i den psykologiske
Beskrivelse og Forstaaelse af en Karakter altid er Noget,
som undslipper os, fordi det er umuligt at gribe alle
Nuancer og fremdrage de individuelle Egenskaber i deres
hele Mangfoldighed, saaledes vil der ogsaa ved Bestemmelsen
af den individuelle Opgave være individuelle
% --- p. 56 ---
\setcounter{footnote}{0}%
\opage{56}Nuancer for hvert enkelt Menneske, som ikke kunne
fremstilles udtømmende i nogen Formel. Eller rettere: i
den samme Formel maa der for hvert enkelt Individs
Vedkommende indsættes særegne Koefficienter. De juridiske
Love skære Alle over én Kam, og de gængse
etiske Love gøre det ogsaa i alt for høj Grad. En fin
Sans for de individuelle Forskelligheder i kvalitativ og
kvantitativ Henseende er endnu overordentligt sjelden.

Den samme Tankegang, der lod den etiske Vurdering
gaa tilbage fra Handlingen til Motivet, maa konsekvent
føre et Skridt videre og tage Hensyn til de Anlæg og
Dispositioner, der betinge Motivernes Opstaaen. Muligheden
for et Motivs Opstaaen er ikke den samme hos
alle Individer. Kun en absolut indeterministisk Opfattelse
vilde kunne antage, at (som en tysk Jurist\footnote[1]{Citeret i Vierteljahrsschrift für wissenschaftliche Philosophie. VI. p. 201.} har udtalt
det) »den moralske Evne til at hæmme lovstridige Drifter
er en uforanderlig Størrelse, som normalt er at statuere
hos alle Mennesker uden Hensyn til individuelle Anlæg«.
En saadan Antagelse strider dog i den Grad mod al
psykologisk Erfaring, at den vistnok vil blive sjeldnere
og sjeldnere i videnskabelige Kredse. --- Og desuden er
det ikke nok med, at et Motiv kan gøre sig gældende
med en vis Styrke. Den Indflydelse, det kan faa paa
Villien, vil bero paa de andre Motivers Styrke i Forhold
til det. Saaledes kunne, som allerede Butler\footnote[2]{Sermon, XII.} har bemærket,
to Mennesker have nøjagtigt samme Grad af
Medlidenhed; men hos den Ene rører der sig en meget
stærk Ærgerrighed eller Hævnlyst, hos den Anden ere
disse Følelser meget svage; de ville da være stillede
højst forskelligt i etisk Henseende. Det gaar her som
% --- p. 57 ---
\setcounter{footnote}{0}%
\opage{57}med Vægtskaale: hvilken Skaal der skal synke, beror
ikke paa de Lodder, der lægges paa den enkelte Skaal,
men paa deres Forhold til de Lodder, der lægges paa
den anden Skaal.

2. Den første Tænker, som har havt sin Opmærksomhed
henvendt paa det her fremdragne etiske Problem,
er \emph{Aristoteles}. Idet han fremstiller sin bekendte Lære
om Dyden som den rette Midte, indskærper han udtrykkeligt,
at objektivt set (κατ’ αὐτὸ τὸ πρᾶγμα) er Midten
den samme for Alle, idet der kan angives to Yderligheder,
som det gælder at holde sig lige meget borte fra,
men sukjektivt set (πρὸς ἡμᾶς) forholder det sig paa Grund % PRINTED AS IS (sukjektivt: for subjektivt; S1 and the two OCRs read k)
af de individuelle Dispositioner anderledes. Lige saa lidt
som et og samme Kvantum Føde passer for den øvede
Atlet og for den første Begynder i Gymnastiken, lige
saa lidt vil for Alle det indre Midtpunkt for Driftslivet
være beliggende paa samme Sted (τοῦτο δ’οὐχ ἓν οὐδὲ
ταὐτὸν πᾶσιν)\footnote[1]{Eth. Nic. II, 5. p. 1106 a, 32.}. Hvad der hos den ene maa kaldes
Tapperhed, idet det hos ham forudsætter Betvingelsen af
en stærk Frygtagtighed, det vil kunne øves aldeles uden
Villiesanspændelse og Selvovervindelse af den Anden.
Denne Betragtningsmaade gennemfører Aristoteles paa en
interessant Maade med Hensyn til de enkelte Dyder. ---
Denne aristoteliske Lære, som vidner om et overraskende
fint Blik for den individuelle Personligheds Betydning og
Ret i etisk Henseende, staar, som jeg et andet Sted\footnote[2]{„Etik“, p. 133 f.}
har gjort opmærksom paa, i en mærkelig Disharmoni til
den aldeles fremherskende Plads, som Aristoteles giver
Staten i Forhold til Individet. Aristoteles har ikke gjort
noget Forsøg paa at vise, hvorledes hin tilsyneladende
ultra-individuelle Lære kan forliges med den objektive,
% --- p. 58 ---
\setcounter{footnote}{0}%
\opage{58}sociale Betragtningsmaade. Socialt fordres der dog --- saa
synes det i det Mindste --- ofte med Nødvendighed Mere
af det enkelte Menneske, end han kan yde efter den véd
\emph{hans} Natur bestemte »rette Midte«. I Forholdet mellem
Individ og Samfund kommer der en Række Problemer
frem, som de antike Tænkere endnu ikke havde noget
klart Blik for. Alligevel staar hin aristoteliske Lære som
en af de mest geniale Synsmaader, der nogensinde er
gjort gældende i Etiken\footnote[1]{I nyere Tid er Tanken om den etiske Lovs Individualisering gjort gældende af \emph{Jakobi}, \emph{Schleiermacher} og især af \emph{Beneke} (Physik der Sitten. 1822), samt i vor egen Literatur af A. S. \emph{Ørsted} i hans smukke Afhandling „om Grænserne mellem Teori og Praxis i Sædelæren“ (1811), som er aftrykt i første Bind af „Eunomia“.}.

Man kunde maaske mene, at det er et i og for sig
uløseligt Problem, som Aristoteles uden at vide det har
stillet ved sin Lære om den individuelt bestemte »rette
Midte«. Det kunde synes, at et Valg maatte træffes
mellem to lige mislige Udveje. Den ene Udvej vilde være
en etisk Subjektivisme, i Følge hvilken enhver almengyldig
etisk Lov falder bort, da Loven bliver aldeles
forskellig for hvert enkelt Individ. Den anden Udvej
vilde være en uforenelig Strid mellem det enkelte Menneskes
individuelle Etik og det sociale Menneskelivs etiske
Krav. --- Den Opfattelse af den etiske Lov, som i det 
Foregaaende er gjort gældende, hjælper os dog ud over
dette Dilemma.

Subjektivismen afvises ved, at trods de forskellige
Former, den etiske Lov maa antage i og for de forskellige
Individer, kan dog \emph{Retningen} for deres Stræben
være den samme. De etiske Love formulere, som vi allerede
fra et andet Synspunkt have set, Tendentser eller
% --- p. 59 ---
\setcounter{footnote}{0}%
\opage{59}Retninger for Livsførelsen og Livsordningen. Naar flere
Skibe krydse op mod Vinden, opnaa de --- paa Grund
af deres forskellige Bygning --- med lige stor Anstrengelse
forskellige Resultater; men derfor kunne de alle bevare
Retningen; denne ligger i Fortsættelsen af den Kurs, de
holde. Den objektive Betragtningsmaade behøver altsaa
ikke at falde bort, fordi man hævder den etiske Lovs Individualisering.
De individuelle Forskelligheder i Anlæg
og Drifter maa Etiken tage som faktiske Udgangspunkter,
der ikke kunne udslettes, men af hvilke det gælder at
bringe det mest Mulige ud. De medføre ganske vist en
Skranke for Etiken som Videnskab, men en Skranke,
der ikke blot vidner om vor Tænknings Begrænsning,
men først og fremmest om Livets rige Mangfoldighed,
som ikke kan indpasses under bestemte Rubriker, skønt
det, der rører sig i de mange individuelle Existentser,
dog kan være rettet mod det højeste Maal for al menneskelig
Handlen. Dette Maal er tillige saa ophøjet, at
de individuelle Forskelligheders Betydning for den ideale
Betragtningsmaade indskrænkes i høj Grad, om den ikke
ganske forsvinder.

Den anden tilsyneladende Konsekvents bortfalder,
naar man nøjere gennemtænker Forholdet mellem Individ
og Samfund, og naar man tillige opfatter de etiske Loves
egentlige Betydning paa rette Maade. --- Det enkelte
Individ er fra først til sidst Led af Samfundet, og \emph{Samfundets
Liv er det Liv, der rører sig i alle de
enkelte Individer}. Samfundets Maal er derfor ikke
naat, saa længe dets Fordringer til det enkelte Individ
medføre en ubodelig Disharmoni i dettes Indre. Naar
Samfundets Villie ikke kan gennemføres, uden at det
enkelte Individ kues eller knækkes, saa tyder det paa
en Ufuldkommenhed i selve Samfundets Ordning. Idealet
er først naat, naar det Arbejde, Individet øver i Sam%
% --- p. 60 ---
\setcounter{footnote}{0}%
\opage{60}fundets Tjeneste, tillige tjener til Udfoldelse af dets egne
Evner og Drifter paa fri og harmonisk Maade. Den
største Velfærd er kun til Stede, hvor hvert enkelt Individ
udvikler sig paa selvstændig Vis og netop derved
paa bedste Maade fremhjælper en lignende Udvikling hos
de andre Individer. Derved kommer det paa én Gang til
at staa som Maal og Middel. Det Afgørende er, at der
i det enkelte Individs Indre øves et Arbejde i \emph{Retning}
af det etiske Formaal; \emph{hvor meget} Arbejde der kan
øves, vil afhænge af Kræfterne og variere med disse.
Enhver Lov, der vil fastsætte ikke blot Arbejdets Retning,
men ogsaa dets Kvantum (»Mandevilliens qvantum
satis«, for at tale med Henrik Ibsen), kan kun angive et
Gennemsnit eller fastslaa et Minimum. Saaledes gaar det
med de etiske Love i deres historiske Form. Den Maalestok,
den Enkelte til given Tid og i et givet Folk anvender
paa sig selv, vil stedse bære Præget af Tid,
Stand og Race; han ser sig selv, som han tænker sig,
at en upartisk Tilskuer blandt de Mennesker, han kender,
vilde se paa hans Færd. Og selv om de Fordringer, han
herved føler stillede til sig, mere fyldestgøres af ham ved
den almindelige Retning, det besjælende Sindelag i hans
Villen, end ved en Række adækvate Handlinger, saa
har det dog sin \emph{pædagogiske} Betydning, at Fordringerne
fremtræde for ham ikke blot med Idealitetens men ogsaa
med Objektivitetens Præg. Des snarere og lettere vil
Driften vækkes til at komme ud over den hidtilværende
Tilstand. Kontrasten mellem Idealet og den sociale Fordring
paa den ene Side, og Individets faktiske Villen paa
den anden Side, kan ofte være et nødvendigt Gennemgangsled
for, at den etiske Bestræbelse, saaledes som
den i Følge Individets egen Natur er mulig, kan komme
i Stand. En Gnist er nødvendig for at udløse den bundne
Kraft. Der er Naturer, som først gennem en smertelig
% --- p. 61 ---
\setcounter{footnote}{0}%
\opage{61}Konflikt med de objektive Normer faa deres Kraft tilstrækkeligt
frigjort, for hvem der altsaa ingen anden Vej
gaar til Harmonien end gennem Disharmonien. --- De
to Sætninger: at den etiske Lov skal formulere Livets
Retning, og at den skal virke pædagogisk, hænge nøje
sammen. Begge Fordringer fyldestgøres nemlig ved, at
den etiske Lov opstiller \emph{Typer eller Forbilleder}. Et
Forbillede anviser Retningen og opfordrer samtidigt til
Efterfølgelse. Men der fremtræder dog her en tydelig
Forskel mellem den videnskabelige Etik og den praktisk-populære
Etik, mellem Moralfilosofi og Moralprædiken.
Praktisk har det sin store Betydning, at Maalet sættes
højt for Alle, for at deres Kræfter kunne udløses. Det
synes jo at være nødvendigt at forlange for meget, for
at det Tilstrækkelige kan blive ydet. Den videnskabelige
Etik kan i sin Bestemmelse af den etiske Lovs Begreb
ikke tage Hensyn til den praktiske Teknik. Hvor let
eller vanskeligt det er i Praxis at finde den for hvert
enkelt Individ efter dets særegne Natur gældende Lov,
gør videnskabeligt set Intet til Sagen. --- Paa den anden
Side er der vel heller intet Spørgsmaal om, at jo mere
individualiseret Moralprædikenen bliver, des mere vil den
virke. Moraliseren en bloc hjælper neppe meget. Det er
altsaa --- trods den tilsyneladende Modsætning --- dog et
Ideal, den videnskabelige Etik her fremkalder for den
praktiske Moraliseringskunst.

3. Man kunde sige: »Den etiske Lov er jo en af
Velfærdsprincipet udledet Konsekvents, ved hvilken der
ikke kan rokkes af Hensyn til det enkelte Individs Natur.
Hvad der i Kraft af et almindeligt Princip er ret og
rigtigt, maa være det for \emph{alle} Individer.« Men, som
allerede vist, hører hvert \emph{enkelt} Individ med sine Forudsætninger
og Anlæg dog med til den Kreds af Væsener,
hvis Velfærd skal fremmes. Dersom der stilles en ubøn%
% --- p. 62 ---
\setcounter{footnote}{0}%
\opage{62}hørlig Fordring til det om Noget, som virkeligt ligger
over dets Evne, saa er der et mod selve Velfærdsprincipet
stridende Misforhold mellem Samfund og Individ til
Stede, et Misforhold, hvis Aarsag ikke \emph{behøver} at ligge
\emph{ene} hos Individet. Den Lidelse, Individet kommer til at
udholde paa Grund af Misforholdet mellem dets Natur
og Samfundets Krav, er derfor ikke blot at opfatte som
Straf, men ogsaa som en Skæbne, lige som den Lidelse,
der udspringer af den manglende Overensstemmelse mellem
Individets Natur og de ydre Naturforhold. Naar det
viser sig, at Ændringer i Samfundets Ordning og i de
Livsbetingelser, som derved beredes for de enkelte Individer,
ofte have langt større Indflydelse paa Forbrydelsernes
Talrighed end Reformer af Straffevæsenet\footnote[1]{G. \emph{Tardé}: La criminalité comparée. p. 61: Changez les conditions, s'il se peut, de la société, bien plutôt que son système de pénalité, et sa criminalité se modifiera. --- Om de sociale Aarsager til Forbrydelse se \emph{Mischler} i „Handbuch des Gefängsnisswesens“, herausg. von Holtzendorff und Jagemann II. p. 447. 481 ff.}, saa
er det klart, at man ikke kan stille det historisk givne
Samfund med dets Krav som absolut berettiget paa den
ene Side, og det enkelte Individ med dets bestemte Natur
og Dispositioner som absolut forpligtet paa den anden
Side. Det enkelte Individ er jo overfor Samfundet aldrig
blot Aarsag, men altid tillige Virkning, og --- hvad der
især maa fremhæves --- det er Virkning, før det er Aarsag.
Selv om der i Individets Natur laa ulykkelige, antisociale
Dispositioner, saa vilde disse maaske ved en bedre
Ordning af Samfundet kunne hæmmes, afledes eller omsættes
i uskadelige Former, medens den nuværende Ordning
maaske ligefrem yder dem Betingelser for en frodig
Udvikling. Og jo mere Samfundets Ordning sker efter
% --- p. 63 ---
\setcounter{footnote}{0}%
\opage{63}bevidst Overlæg, des større bliver dets Ansvarlighed
overfor det enkelte Individ.

Dersom denne Opfattelse er rigtig, maa man vente
at finde Spor af den ogsaa i Teologien og Retslæren,
skønt disse gaa ud fra Modsætningen mellem en absolut
objektiv Lov paa den ene Side og Individet, der har at
lyde den, paa den anden Side. Jeg finder i den Rolle,
som Naade og Benaadning spille i Teologien og Retslæren,
et Spor af den Tanke om den sande Lovs Individualisering,
som jeg her har søgt at udvikle. I Naaden og
Benaadningen kundgør der sig nemlig mere eller mindre
tydeligt et Hensyn til den Handlendes Individualitet og
til dennes subjektive Betingelser. Den ortodoxe Lære
antager jo endog, at det naturlige Menneske ikke \emph{kan}
opfylde Loven! Naaden underkender altsaa Loven ved
at hævde, at Menneskets hele Individualitet ikke rammes
af Modsætningen til Lovens Bud. I selve Muligheden af
en Forsoning ligger, at Loven ikke var absolut berettiget
overfor Individet; --- ellers vilde Forsoningen være en
ligefrem Selvmodsigelse. Og naar efter den ortodoxe
Forsoningslære Gud selv, Lovens Giver, bringer det
Offer, ved hvilket Lovens Overtrædelse sones, kunde man
heri finde et symbolsk Udtryk for, at det er Loven, ikke
Mennesket, der har haft den egentlige Uret. --- Efter den
juridiske Opfattelse er Benaadningen ligefrem »Retfærdighedens
Selvkorrigeren«\footnote[1]{\emph{Jhering}: Der Zweck im Recht. I. p. 428.}, eller som en ældre Jurist\footnote[2]{\emph{Anselm v. Feuerbach}: Aktenmässige Darstellung merkwürdiger Verbrechen. Giessen 1828. I. p. 353.}
udførligere udtaler det: »At bøde paa Lovenes Ufuldkommenheder
under deres Anvendelse, med Visdom at
udligne Misforholdet mellem den lovlige Straf og en enkelt
Overtræders Strafbarhed, mellem den uforanderlige, al%
% --- p. 64 ---
\setcounter{footnote}{0}%
\opage{64}mindelige Lovs ubøjelige Strenghed og den individuelle
Skylds mangfoldige Foranderligheder, og saaledes \emph{at forsone
Retfærdigheden med Billigheden}, denne Magt
er med Benaadningsretten lagt i Suverænens Haand«.
Dersom selve Retfærdigheden trænger til at korrigeres,
maa det komme af dens Ufuldkommenhed! --- Ved den
juridisk-teologiske Modsætning mellem Synd og Naade,
Forbrydelse og Benaadning svinges der til begge Sider
af det Midtpunkt, i hvilket den egentlige Sandhed ligger.
Paa Grund af Vanskeligheden ved at ramme selve Midtpunktet,
bevæger man sig først et Stykke ud over det
til den ene Side, dernæst et Stykke til den modsatte
Side. I Stedet for at bevæge sig efter Diagonalen, bevæger
man sig først ad den ene, dernæst ad den anden
af Parallelogrammets to Sider. Og som allerede antydet,
der synes at være Naturer, som kun kunne udvikle sig
gennem saadanne Modsætninger, og som ikke kunne bevæge
sig efter Diagonalen.

I den moderne Udvikling af Straffevæsenet finder
jeg mere eller mindre bevidst den samme Tanke ligge til
Grund. --- Naar det Minimum af ydre Sædelighed, Staten
fordrer af sine Medlemmer, i Virkeligheden stiller forskellige
Fordringer til de Enkelte efter disses forskellige
Natur, saa er det Statens Pligt at tage Hensyn til disse
Forskelligheder og finde Udgangspunkter i Individernes
Natur for en Udvikling henimod det, Loven fordrer.
Alt er ikke gjort med at fastsætte Loven og idømme
Straffen. Den Erkendelse synes da ogsaa nu at være 
vaagnet, at Straffefuldbyrdelsen er af nok saa stor Vigtighed
som Straffelovgivningen og Straffedommene, og at
de Erfaringer, der indsamles under Straffefuldbyrdelsen,
bør blive af væsentlig Betydning for Fremtidens Strafferet.
I Straffefuldbyrdelsen kommer netop det individualiserende Moment frem, som Straffelov og Straffedom kun
% --- p. 65 ---
\setcounter{footnote}{0}%
\opage{65}højst ufuldkomment kunne lade lade komme til dets Ret. % PRINTED AS IS (ufuldkomment: lade lade: doubled (600 ppi))
De Erfaringer, som gøres ved Straffefuldbyrdelsen, ville
derfor mere og mere virke tilbage paa Strafferettens
Ordning. Saaledes har \emph{Franz von Holtzendorff}
nyligt udtalt: »De moderne Straffeteorier maa langt snarere
udledes ad Induktionens Vej og \emph{af Straffefuldbyrdelsens Erfaringer} end omvendt Straffefuldbyrdelsen af aprioriske
Teorier«\footnote[1]{Handbuch des Gefängnisswesens. Hamburg 1888. I. p. 387. --- De følgende Bemærkninger om Fængselsvæsenets Udvikling ere ogsaa hentede fra dette Værk.}. Det anerkendes tillige mere og mere, at
det ikke er Forbrydelserne som objektive Fakta, men Forbrydernes
Personligheder, man i Straffevæsenet har med
at gøre. Man har derfor fordret de uforanderlige Domme
afløste af Domme, som fastsætte en ubestemt Straffetid,
for at der kan tages Hensyn til Individets Karakter, saaledes
som denne viser sig under Straffefuldbyrdelsen. Man
fordrer en bestemt Inddeling af Forbryderne med forskellig
Behandling efter deres forskellige Karakter. Der
gør sig i det Hele i det sidste Aarhundrede (efter John
Howards Tid) ved alle Reformer af Straffevæsenet den
Tanke gældende, at Forbryderen saa vidt muligt skal træde
tilbage i det menneskelige Samfund ikke blot kuet og
tvungen, men virkeligt forbedret. Den Erfaring er bleven
almindeligt anerkendt, at for Mange ligger efter deres
indre og ydre Forhold Forbrydelsen uendeligt nær, medens
den for Andre ligger meget fjernt. Det enkelte
Individs Ret overfor den abstrakte, upersonlige Lov anerkendes
saaledes mere og mere. Individualisering og
Humanisering gaa Haand i Haand og ere i Virkeligheden
 uadskillelige.

I en ideal Stat vilde der kun fordres af Enhver hvad
der svarede til hans Evner og Kræfter. Fordringen vilde
være afpasset efter hans Ejendommelighed, ligesom vi ved
% --- p. 66 ---
\setcounter{footnote}{0}%
\opage{66}Opdragelsen afpasse de Fordringer, vi stille til Barnet, efter
dettes Udviklingstrin, for at Fordringen kan give Anledning
til Udfoldelse af Kræfterne og ikke til Underkuelse
af hele Individets Væsen. De sociale Ulemper, der vilde
hidrøre fra, at de samme Fordringer ikke stilledes til alle
Individer, vilde da være at afhjælpe ved særegne Foranstaltninger.
Saa længe dette Ideal af en gennem alle
sine Institutioner og Love opdragende Stat ikke kan realiseres,
maa Staten gennem Straffefuldbyrdelsen øve en Del
af det Værk, den ikke paa en mere direkte etisk Maade
har formaaet at fuldbyrde. Den virkelige Ordning fører
til at svinge om den sande, ideale Ordning som usynligt
Midtpunkt i haarde, smertelige Modsætninger.

4. Ved at tage Hensyn til den individuelle Faktor
ville de kasuistiske Problemer løses som af sig selv. Rent
objektivt kan det i det enkelte Tilfælde være umuligt at
afgøre, om der snarere skal arbejdes for et større Samfund
eller for et mindre, om Hovedvægten skal lægges
paa Udbredelsen af den vundne Kultur til saa Mange
som muligt eller paa at frembringe ny Kultur, som foreløbigt
kun vil komme nogle Faa til Gode, o. s. v. Men saadanne
Spørgsmaal ville i Regelen finde deres Afgørelse
derved, at Enhver har sine Evner og Anlæg, som betinge
hans Kald. I Begrebet Kald ligger netop tydeligt udtalt,
at et Menneskes Pligt ikke blot bestemmes ved de
almindelige etiske Principer og de givne faktiske Forhold,
men ogsaa ved hans egen individuelle Natur. Begrebet
Kald hviler paa den Tanke, at Mesterskabet viser
sig i Begrænsningen, da al Villen og Handlen maa gaa
ud paa noget Bestemt, og at Begrænsningen væsentligt
beror paa de individuelle Evner og Kræfter. Det er
disse, der »kalde«, idet Kraften har en Tendents til at
vække Driften. --- Tage vi Virksomheden i Kulturens
Tjeneste som Exempel, saa vil der være dem, hvis Evner
% --- p. 67 ---
\setcounter{footnote}{0}%
\opage{67}bedst gøre dem skikkede til at virke i en snever Kreds,
og hvis Virksomhed derfor naturligt bliver indskrænket
til at sørge for deres egen og deres Nærmestes Opholdelse,
saaledes at deres Virken i det Højeste kun indirekte
kan faa Betydning ud over den snevrere Kreds. Andre
vil der være, som egne sig til at virke paa og for Mange,
og som derfor ville være gode Redskaber til at udbrede
den hidtil vundne Kultur videre. Andre igen ville have
Instinkt og Kraft til at slaa ind paa nye Veje, til at
prøve nye Muligheder. Trods Modsigelse og Mangel paa
Forstaaelse fra Mængdens Side øve de dog deres Værk
og pleje de første Spirer til Fremtidens Kultur. Der vil
for den Enkelte, der kender sig selv, her ikke være nogen
Tvivl om, hvor han skal vende sig hen. Men Hovedvanskeligheden
er netop den at kende sig selv, som
allerede \emph{Sokrates}, Etikens Grundlægger, indskærpede.
Den sokratiske Grundtanke viser sig netop ved den Opfattelse,
jeg her søger at gøre gældende, af Loven som
individuel, i sin hele store og dybe Betydning. Selvskuffelse
er saa let mulig, hvor det gælder at udfinde
sin egen Evne og afstikke sit Maal og dermed sin Grænse.
Intet objektivt Kriterium slaar her ganske til, og der
maa stedse voves noget, handles i Tillid til hvad den
personlige Overbevisning forkynder. Ethvert Valg er
her en praktisk Hypotese, hvis Verification det er Livets
Opgave at bringe. Hvis Individet i sin etiske Formel
anbringer urigtige Koefficienter, bliver Opgaven ikke løst.

Dette gælder ogsaa, naar vi i Stedet for Kaldsvirksomheden
tage \emph{Selvbeherskelsen} som Exempel til at
belyse Forholdsloven i Etiken.

Selvbeherskelse paa enkelte Omraader kan forekomme
sammen med fuldstændig Mangel paa Selvbeherskelse paa
andre Omraader. Indianerne f. Ex. ere i Stand til at
udholde de forfærdeligste Pinsler, naar de falde 5\% i Fjende%
% --- p. 68 ---
\setcounter{footnote}{0}%
\opage{68}haand, eller naar de ved at paalægge sig selv de mest
udsøgte Lidelser ville forsone og ære deres Guder. Ogsaa
Sygdom bære de uden Klage, og de fordre af deres
Kvinder, at de skulle føde uden Smerteskrig. I deres
ydre Optræden bevare de stedse den største Ro, Anstand
og Alvor, selv om de heftigste Lidenskaber rase i deres
Indre. Men overfor Nydelser som Spise, Drikke og Spil
formaa de ikke at styre sig. Det var med stor Anstrengelse,
at nogle Høvdinge fik sat igennem, at Drukkenskab
skulde betragtes som en Last. --- Dette viser, at
Evnen til Selvbeherskelse ikke er en konstant Faktor,
uafhængig af hvad der ellers rører sig i Mennesket.

Endnu klarere viser dette sig ved Selvbeherskelse
paa Kønsinstinktets Omraade. De forskellige Individer
ere her højst forskelligt udrustede. Der gives dem, hos
hvem Instinktet virker i Harmoni med sympatiske Følelser,
saa at det ikke kræver Tilfredsstillelse, uden hvor det
danner det naturlige Grundlag for inderlig Hengivelse til
et andet Individ, der omfattes med Glæde og Beundring.
Dette er af saa meget des større Betydning, som Tilfredsstillelsen
af Kønsinstinktet hovedsageligt af den Grund
falder ind under etisk Betragtning, at den ikke er mulig
uden ved Hjælp af et andet Individ. Afset fra dette
Forhold og det deraf udspringende Ansvar vilde Kønsinstinktet
kun være at bedømme fra et rent individualistisk
eller egoistisk Standpunkt, altsaa efter hygiejniske
Hensyn eller efter Hensynet til den harmoniske Udvikling
af Individets sjælelige og legemlige Evner i det Hele.
Det store Problem, den etiske Vanskelighed fremkommer
netop ved, hvorledes dette mægtige Instinkt, der sætter
saa stærke Følelser i Bevægelse og er saa betydningsfuldt
for selve det enkelte Individs Velfærd, skal kunne
tilfredsstilles saaledes, at det ikke volder fysisk og moralsk
Lidelse for andre Individer og griber forstyrrende ind i
% --- p. 69 ---
\setcounter{footnote}{0}%
\opage{69}den rette sociale Ordning. --- Foruden de før omtalte
lykkeligt anlagte Individer gives der andre, hos hvem
Instinktet rører sig i hele sin elementære Brutalitet, som
en rent fysisk Trang, ved hvis Tilfredsstillelse Genstanden
selv kan være ligegyldig. --- Hos atter Andre er Instinktet
forholdsvis stærkt, men de formaa dog at aflede Opmærksomheden
fra det ved ideelle Sysler eller andet Arbejde.
--- Hos atter Andre kan det ikke beherskes paa denne
Maade. Der findes i det Hele her utallige Gradationer lige
til den »sygeligt forhøjede Kønsdrift«, som i Følge Krafft-\emph{Ebing}
hyppigt findes hos Individer med en »neuropatisk
Konstitution«, Individer, der »en stor Del af deres Liv
lide under Byrden af denne konstitutionelle Anomali i
deres Driftliv«\footnote[1]{Psychopathia sexualis. 2. Aufl. p. 34.}. --- De to Køn ere i Regelen forskelligt
disponerede i denne Henseende. Harmonien mellem Kønsinstinktet
og den ideelle Hengivelse til et andet Individ
er sikkert langt hyppigere hos Kvinden end hos Manden,
hvilket turde hænge sammen med, at Hannen i hele Dyreverdenen
er den mest aktive og agressive Part ved Forplantningen,
hvorfor Instinktet hos den har en saa elementær
og hensynsløs Styrke, som det kun sjældent har hos
Hunnen. --- Men ogsaa hos Mændene varierer Instinktets
Styrke og Forhold til andre Motiver som sagt i det
Uendelige. Under den Strid, som hyppigt føres om sexuel
Sædelighed, overser man i Regelen disse talløse individuelle
Forskelligheder, som dog sikkert er det aller vigtigste
Element i den hele Sag. Man opstiller fra de modmodsatte % PRINTED AS IS (modmodsatte: mod-|modsatte at the line turn (600 ppi))
Sider lige absolute Paastande. Fra den ene
Side formulerer man strenge, almengyldige Bud uden at
indse, at man derved lægger en højst forskellig Byrde
paa de forskellige Skuldre. Fra den anden Side postulerer
man en fysiologisk Nødvendighed, som Enhver skal
% --- p. 70 ---
\setcounter{footnote}{0}%
\opage{70}ligge under for. Dette er et Fatum, som ikke existerer.
Kønsinstinktet fremtræder under andre Forhold hos Mennesket
end hos Dyret. Naar det vaagner hos Mennesket,
har Evnen til Selvbeherskelse allerede paa andre Omraader
udviklet sig. Der er sat Forestillinger og Følelser
i Bevægelse, som ikke tillade Instinktets direkte Herredømme,
og som, naar Opmærksomheden og Interessen for
ideel Selvopholdelse er vakt, kunne gøre det muligt at
dæmpe det og bringe dets Ytringer i Harmoni med Livets
øvrige Krav. Kun maa det ikke glemmes, at Selvbeherskelse
som rent negativ Dyd er en psykologisk Umulighed.
Det er en Sætning, som altfor ofte glemmes i Etiken, at
en Drift kun kan hæmmes af en anden Drift. Det kommer
derfor ved Selvbeherskelsen ikke blot an paa, hvor stærk
den Tilbøjelighed er, der skal kues, men ogsaa paa, om
der er Mulighed til Stede for andre Tilbøjeligheder, der
kunne lægge Beslag paa den sjælelige Energi. Det gælder
om, at der er en stor Kærlighed, en stor Begejstring,
som kan fylde Sjælen. Forholdsloven paa Etikens Omraade
viser sig her at være en simpel Følge af Loven
om Energiens Bestaaen\footnote[1]{Man vilde her ogsaa kunne benytte sig af Inertiens Lov. \emph{Spinoza} beviser sin Sætning: affectus coërceri nec tolli potest, nisi per affectum contrarium et fortiorem (Eth. IV, 7) ud fra Inertiens Lov.}. Den Energi, vor Bevidsthed
eller vor Hjerne raader over, er til enhver Tid begrænset;
jo mere af den, der forbruges til én Drifts Tilfredsstillelse,
des mindre bliver der tilbage til de andre Drifter. Tilintetgørelse
er umulig, men Omsætning til andre Former
kan være mulig. Selvbeherskelsens store Heroer have
besiddet en stor Begejstring for andre Sider af Livet end
dem, de vendte sig fra. Den store Selvbeherskelse, Indianerne
vise overfor Saar og Smerte i Krig, Fangenskab
eller religiøs Dyrkelse, betinges ved deres Æresbegreber,
% --- p. 71 ---
\setcounter{footnote}{0}%
\opage{71}deres Stammefølelse, deres religiøse Tro. Overfor sanselige
Udskejelser havde deres Livsforhold ikke givet Anledning
til opvejende Motivers Opstaaen. Som \emph{Leslie
Stephen} har bemærket, udvikler der sig hos Mennesket,
efterhaanden som det skrider frem i intellektuel, æstetisk
og sympatisk Henseende, en højere Sensibilitet, der medfører
en Følelse af Væmmelse og Foragt for sanselige
Udskejelser afset fra disses beregnelige Følger, en Sensibilitet,
der kan virke lige saa direkte og imperativisk
som nogen af de elementære Sanseevner\footnote[1]{Science of Ethics. p. 102. 195. 200. --- Se ovenfor p. 23.}.

Selvbeherskelsens Kamp staar paa, indtil den nye
Anvendelse af Energien har faaet fuldstændig Overvægt
over den tidligere. Denne Kamp bliver af højst forskellig
Art hos de forskellige Individer, efter Forholdet mellem
de oprindelige eller erhvervede Tilbøjeligheder. Hos
Nogle foregaar Omsætningen endog uden Kamp; men der
er dem, som kæmpe haardt uden at naa fast Grund under
Fødderne. Det Villiesarbejde, som foregaar i disse sidste,
er langt alvorligere end det, der øves af hine, skøndt det
ydre Resultat bliver langt ringere. Naar Jesus hellere
omgikkes med Skøger end med Farisæere, var det maaske,
fordi han hos dem fandt en indre Stræben og Higen, der
vel ikke kunde overvinde de modstaaende Skranker, men
dog var i nærmere Slægtskab med den ideale etiske Lov
end Farisæernes selvtilfredse Korrekthed. ---
--- Forholdsloven kaster Lys over mange etiske Spørgsmaal,
der ellers vilde være uløselige. Den lærer os at
forene Etikens Idealitet med den realistiske Anerkendelse
af Livets Mangfoldighed, idet den viser os, at den egentlige
etiske Udvikling foregaar ud fra en Mængde forskellige
Udgangspunkter, fra meget forskellige Niveauer
henimod et og samme Maal.

% --- p. 72 ---
\setcounter{footnote}{0}%
\opage{72}5. Hvor den af de almindelige Principer og den
givne Situation udspringende Opgave ganske stemmer
med Individets egen Evne og Drift, have vi de lykkeligste
--- og naar Opgaven er stor --- de største Mennesker.
Det Karakteristiske for disse er den indre Sikkerhed
og Harmoni, med hvilke de --- trods alle indre og ydre
Kampe --- gaa frem mod deres Maal og udfylde deres
Plads. Her er den af \emph{Herbert Spencer} saa kaldte
»organic morality« til Stede. At der gives saadanne
Naturer er et Faktum, som ikke kan benegtes, hvorledes
man saa vil forklare det. Det er dem, der afgive Forbilleder
for Andre. Man maa herved vel fastholde, at
Opgaven ikke blot stilles Individet ud fra uforanderlige,
én Gang for alle givne Situationer, men at Individet, i
Kraft af sin Evne og Drift og i Henhold til de højeste
Principer ogsaa netop kan stille sig selv den Opgave at
forandre den foreliggende Situation, hvortil denne saa
naturligvis selv maa indeholde en Mulighed. Hos de
store Banebrydere udvikle sig Evner og Drifter, som
ikke kunne finde deres Tilfredsstillelse under de givne
Forhold, og som derfor selv ud af disse frembringe nye
Forhold. Saadanne Individer opstaa sandsynligvis ved
en af de Variationer, der i Følge Darwins Hypotese ere
Forudsætningen for, at en ny Livsform kan danne sig.
Deres Ret falder sammen med deres Evne. Ved selve
Udfoldelsen af hvad der ligger i deres Natur afgive de
Forbilleder for en ny Retning paa et af Menneskelivets
Omraader og blive derved Midler for Slægtens Fremgang.
Natur og Kunst, Natur og Pligt virke her umiddelbart
sammen. De bære Loven i sig selv i højere Forstand
end det kan siges om dem, i hvilken ingen nye Muligheder
for en højere Udvikling ligge. Dog kan der ogsaa
hos dem, hvis Opgaver væsentligt bestemmes ved Forholdene
som de ere givne, findes »organisk Moralitet«,
% --- p. 73 ---
\setcounter{footnote}{0}%
\opage{73}dersom deres Natur umiddelbart fører dem baade til den
Virksomhed og til den Begrænsning og Selvbeherskelse,
som Forholdene kræve.

Den fuldstændige Modsætning til saadanne Naturer
vilde være til Stede, dersom der skulde gives Individer,
hos hvem aldeles ingen, end ikke den svageste Bevægelse
eller Stræben i Retning af det etiske Ideal kunde udløses.
Alle Koefficienterne i den etiske Formel vilde da være
Nul. Ingen nok saa pædagogisk formuleret Fordring eller
Appel vilde kunne virke som stimulus. Vi komme her
blandt Andet til det Spørgsmaal, om der gives uforbederlige
Forbrydere, et Spørgsmaal, som i den nyeste Tid
med stor Bestemthed er blevet besvaret bekræftende af
\emph{den italienske kriminalpsykologiske og retsfilosofiske
Skole}.

Denne Skole har store Fortjenester af Videnskaben.
Den sædvanlige etiske og juridiske Betragtningsmaade
bliver staaende ved Handlingen og de nærmest liggende
Motiver og gaar ikke tilbage de Aarsager, som have % PRINTED AS IS (tilbage: gaar ikke tilbage de: til missing (600 ppi))
frembragt den Handlendes hele Karakter. Den tidligere
herskende Indeterminisme har begunstiget denne Opfattelse.
Den kriminelle Antropologi og Psykologi, som
skylder den italienske Skole saa meget, arbejder paa at
raade Bod paa denne Ensidighed. Den søger at konstatere
Forbryderens legemlige og aandelige Ejendommeligheder
og at paavise Aarsagerne til disse. Den har
paa Grundlag af disse Undersøgelser opstillet Forbrydertypen
(homo delinquens) som en særegen Variation af den
almindelige Mennesketype. Blandt de legemlige Særegenheder
for denne Type nævner \emph{Lombroso} særligt saadanne
Egenskaber ved Kranie, Hjerne og Fysiognomi,
der maa betragtes som Afvigelser fra det \emph{civiliserede}
Menneskes Type (\emph{regressive} Afvigelser), bl. a. ringe
Kraniekapacitet, lille og snever Pande. Blandt de sjæle%
% --- p. 74 ---
\setcounter{footnote}{0}%
\opage{74}lige Særegenheder fremhæver han en mærkværdig Uimodtagelighed
for Smerte, med hvilken vistnok en ikke mindre
mærkelig Mangel af Medlidenhed staaer i naturlig Forbindelse.
De vasomotoriske Reflexer ere i det Hele svage,
og det er sjeldent, at de egentlige Forbrydernaturer kunne
rødme. Den moralske Følelse mangler, hvad der (som
allerede \emph{Despine} havde udtalt) viser sig i Angerens
store Sjeldenhed. Hele denne Type er nu i Følge Lombroso
og den italienske Skole at forklare ved Atavisme.
Den opstaar nu til Dags ved, at dybe Lag fra Menneskeslægtens
Forhistorie skyde sig op igen midt under Civilisationen.
Paa anden Maade mener man ikke at kunne
forklare det Umotiverede og Hensynsløse ved de egentlige
Forbryderes Færd, ligesom ogsaa Tilbagefaldenes Hyppighed
og de bedst udtænkte Straffesystemers ringe Virkning
forstaas, naar det er Kræfter fra Urtiden, som røre
sig paa ny.

Den moralske Sans --- siger en fremragende italiensk
Forsker, Garofalo (i sin »Crimonologie«) --- er Aarhundreders % PRINTED AS IS (Crimonologie: for Criminologie (600 ppi))
Værk. Hvis nu --- som Erfaringen viser ---
enkelte Individer ikke ere blevne delagtige i dette Værk,
men lide af en medfødt moralsk Insensibilitet, hvilket
viser sig i, at Medlidenheden ikke rører sig hos dem,
end ikke i dens mest elementære Former --- hvorledes
kan man saa tænke sig, at nogensomhelst Paavirkning
og Opdragelse i Løbet af det enkelte Individs Livstid
skulde kunne udrette, hvad der for den øvrige Slægt har
kostet stedse gentagne Erfaringer i Løbet af Aartusinder?
Mennesket er godt af Instinkt, ikke af Reflexion; naar
da Instinktet mangler, hvilken anden Kraft skulde saa
kunne raade Bod derpaa? Individer, der ikke (som
Lejlighedsforbrydere) forsynde sig paa Grund af forbigaaende
Anledninger, men som atter og atter vende tilbage
til deres forbryderiske Færd, eller som ved Hand%
% --- p. 75 ---
\setcounter{footnote}{0}%
\opage{75}lingens forholdsvis ringe Motiver og den oprørende Maade,
paa hvilken den udføres, aabne Indblik i et indgroet antisocialt
Sindelag, --- saadanne Individer er det umuligt
for et civiliseret Samfund at assimilere sig. De vilde
kunne passe i langt tidligere Civilisationer, og nu om
Stunder maaske i Dahomey eller paa Fidschiøerne; men
det er haabløst at ville gøre dem til Elementer i vort
nuværende Samfundsliv. Overfor dem kan Samfundet da
ingen Pligt have. Derimod har det overfor sig selv den
Pligt at eliminere saadanne »Molekuler«, som ville skade
Livet i den øvrige Organisme. Det er en individualistisk
Ensidighed at ville hævde, at forbryderiske Individer af
den beskrevne Art skulde have Rettigheder overfor Samfundet.
Den sociale Nødvendighed er ubetinget, og den
sociale Reaktion taaler ingen Begrænsning.

Saa vidt \emph{Lombroso} og Garofalo. --- Den af disse
Forskere hævdede Anskuelse interesserer os her, for saa
vidt som den antager en absolut moralsk Insensibilitet
hos nogle Individer, som derved vilde udelukkes fra i
nogen som helst Grad at høre til den Mangfoldighed af
Udgangspunkter, fra hvilke den menneskelige Udvikling
bevæger sig hen imod sit Ideal. Ikke blot den individuelle
Relativitet, men ogsaa den psykologiske Relativitet
vilde falde bort. Vi vilde i de egentlige Forbrydernaturers
Sindelag og Tankegang have givet et System,
som var fundamentalt modsat det etiske System, paa hvilket
den menneskelige Udvikling, bevidst eller ubevidst er
bygget. Og der skulde aldeles ingen Mulighed være for
at hint System kunde overvindes af dette uden ved en
Udryddelseskrig. Førend vi gaa ind paa denne Antagelse,
maa vi underkaste den en nærmere Drøftelse.

a. Det er altid usikkert at slutte fra enkelte Handlinger
til et Menneskes inderste Karakter og de Muligheder,
den indeholder. Et Menneskes Karakter og dens
% --- p. 76 ---
\setcounter{footnote}{0}%
\opage{76}Muligheder udtømmes ikke ved den enkelte Handling. I
Følge Garofalo behøver Forbrydelsen ikke en Gang at
være særligt stor, og dog skal Undersøgelsen af Forbryderens
Karakter ved Hjælp af de Oplysninger, der kunne
skaffes til Veje, kunne være tilstrækkelig for en antropologisk
og psykologisk uddannet Dommer til at afgøre,
om Individet hører til dem, der kunne lempe sig efter
det civiliserede Samfunds Krav, eller om det maa »elimineres«.
Ja, end ikke Handling behøves; et Forsøg skal
kunne være tilstrækkeligt!

Lige saa usikker er den Slutning, der drages fra
Fraværelsen af Anger. Den er især konstateret hos
dødsdømte Forbrydere eller hos Fanger som dem, Dostojevski
maatte dele Fængsel med i Sibirien, og som just
ikke vare Genstand for nogen virkeligt pædagogisk Behandling.
Fordi Angeren ikke kan vækkes i den korte
Tid mellem Forbrydelsens Udøvelse og Dødsdommens
Udførelse, og fordi den ikke kunde vækkes hos Forbrydere,
der underkastedes den Behandlingsmaade, som indtil
for kort Tid siden var den sædvanlige i Fængslerne,
deraf kan man dog ikke slutte, at den \emph{under ingen
Omstændigheder} vilde kunne vækkes. --- Tilbagefaldenes
Hyppighed tyder ligeledes paa indgroede, men
ikke nødvendigvis paa uforbederlige Tilbøjeligheder.

Garofalo vil ved Frihedsstraffene ikke have nogen
bestemt Tid fastsat ved Dommen, blandt Andet »fordi
man ikke kan maale den Andel i Forbrydelsen, som de
omgivende Forholds Indflydelse havde, og den, som selve
Individets Anomali havde«. Først efter en længere Iagttagelsesperiode,
i hvilken man kan lære Individets Karakter
at kende, skal Grænsen for Straffetiden kunne
fastsættes\footnote[1]{La Criminologie. p. 395.}. Men hvorledes kan han saa i nogle Tilfælde
% --- p. 77 ---
\setcounter{footnote}{0}%
\opage{77}fordre absolut Elimination? Den Usikkerhed i Kendskabet
til Forbryderens Karakter, som gør, at Frihedsstraffens
Længde ikke forud skal kunne fastsættes, maa
ogsaa gøre det umuligt at fastslaa absolut Elimination
som eneste Udvej! 

b. Naar de uforbederlige antisociale Anlæg forklares
ved Atavisme, saa rejser der sig et nyt Problem, nemlig
hvorledes saa udprægede Ejendommeligheder, der ikke
blot angaa en enkelt Side i Individets Karakter, men
hele Karakteren, have kunnet holde sig inden for Slægten
ved Siden af Ejendommeligheder af helt modsat Art.
Hvorledes har den Traad, gennem hvilken det enkelte
Individ hænger sammen med Urslægterne, kunnet undgaa
at blive spunden sammen med den Traad, ved hvilken
det hænger sammen med de nærmeste Slægtled og med
de Institutioner og Traditioner, som disse have frembragt
og overleveret? Med andre Ord: selv om Forklaringen af
Forbrydertypen ved Atavisme er rigtig\footnote[1]{G. \emph{Tarde} har i sin „Criminalité comparée“ kritiseret denne Forklaring og kommer til det Resultat, at selv om man maa anerkende utvivlsomme fysiologiske og anatomiske Analogier mellem de fødte Forbrydere og de Vilde, og selv om hine --- som alle Monstra --- frembyde regressive Træk, saa findes de dog hos dem \emph{kombinerede paa en særegen Maade}, og man maa ikke bedømme vore Forfædre efter denne Prøve. (p. 46). Men saa gælder det at forklare „Kombinationen“ --- og den kan dog ikke være uafhængig af den sociale Udvikling.}, saa \emph{er Atavismen
intet Bevis paa}, at \emph{Samfundet ingen Skyld
og derfor heller ingen Pligt har over for Individet}.
Ti naar de antisociale Spirer have kunnet holde
sig saa længe i latent Form inden for Slægten og ikke
ere blevne omdannede, maa det dog komme af Mangler
og Ufuldkommenheder i Slægtens Liv. Samfundstilstandenes
og i det Hele det etiske Livs Fuldkommenhed maa
dog maales, ikke blot efter hvilke positive Fremskridt
% --- p. 78 ---
\setcounter{footnote}{0}%
\opage{78}der gøres, men ogsaa efter hvilke uheldige Muligheder
der fortrænges og ophæves. Jo fuldkomnere en Samfundsorden
er, des mere virker den paa de enkelte Individers
inderste Væsen. Atavismen er altsaa selv et Tegn paa
social Ufuldkommenhed og retfærdiggør ikke, at Samfundet
stilles overfor det forbryderiske Individ som den absolute
Ret overfor den absolute Retsløshed. --- Mærkeligt nok
kommer den italienske Skole, idet den stiller de forbryderiske
Individer uden for Samfundet baade i fysisk og
etisk Henseende, til samme Resultat som Indeterminismen
og Individualismen, skønt disse ellers ere dens Antipoder.

c. Naar Garofalo, lige som tidligere \emph{Despine}
(og endnu tidligere \emph{Anselm von Feuerbach}), lærer,
at den moralske Insensibilitet kan forblive latent, indtil
ydre Anledninger røbe den, --- naar han lærer, at selv
om Fattigdom ikke er den egentlige Aarsag til Forbrydelse,
saa vilde dog mange Mennesker være vedblevne
at være latente Forbrydere, hvis de ikke havde været
fattige, og hvis de havde levet i et andet milieu, --- saa
synes det klart, at Samfundet maa være en medvirkende
Aarsag, da dets Ordning altid mere eller mindre betinger
det milieu, i hvilket de enkelte Individer udvikle sig.
Man kan skelne nok saa skarpt mellem den blotte Anledning
og den egentlige Aarsag; denne Distinktion vil altid
være mere eller mindre kunstig, idet \emph{den fuldstændige
Aarsag bestaar baade af de ydre Anledninger
og de indre Anlæg}. Der er saa meget des mere Grund
til at fremhæve dette, som det jo er »Anledningerne«, der
bevirke den uhyre Forskel mellem to latente Forbrydere,
at den Ene bliver erklæret for inassimilabel, medens den
Anden uforstyrret fortsætter sit Liv som Samfundsmolekule!


d. \emph{Nedarvede Anlæg have stedse en vis
Ubestemthed}, der gør det muligt, at de ved Opdragelse
% --- p. 79 ---
\setcounter{footnote}{0}%
\opage{79}og ydre Forhold kunne udvikles i forskellig Retning.
Hvad der nedarves, er kun Dispositioner. Der kan ikke
bevises at være en fatalistisk Naturbestemmelse til Stede,
som skulde gøre alle Erfaringer, Oplevelser og Anstrængelser
betydningsløse. Selv om en god Opdragelse og
lykkelige Forhold ikke skulde formaa at kvæle den onde
Spire, men i det Højeste kunde bevirke, at den udfolder
sig uden sociale Ulemper, saa vil dog en slet Opdragelse
gøre overordentligt Meget til at udvikle den onde Spire.
Man vil derfor ikke kunne drage nogen absolut Grænse
mellem hvad der udspringer af Individet i og for sig, og
hvad der skyldes de Forhold, i hvilke det stedes, og som
for en stor Del bero paa Samfundets Ordning. --- Meget
ofte er der ingen absolut Egoisme til Stede. Hos Forbrydernaturer
kunne Familiefølelserne ofte være meget
stærke, saa at de forbryderiske Instinkter kun ytre sig
over for Fremmede. Her kunde da maaske ligge en
Mulighed for Vækkelse af videre gaaende Sympati, lige
som det er fra Familiefølelserne, at Menneskekærligheden
fra først af har udviklet sig til at omfatte større Kredse.
Et andet Udgangspunkt kunde søges i Forbryderens Selvfølelse
og Trang til at gøre sin Person gældende. Dostojevski
bemærker, at den værste Straf for Forbryderen
vilde være \emph{unyttigt} Arbeide. Dette hænger sikkert
sammen med, at unyttigt Arbejde, som det at øse i Danaidernes Kar, lader Personligheden staa som et Intet, der
hverken har Betydning for sig selv eller Andre. Denne
sidste Rest af Selvfølelse vilde uden Paavirkning ude fra
ikke være nok til at overvinde Ladheden, der er en hyppig
Del af Forbrydernaturen; men den kan afgive et Udgangspunkt for Opdragelse, og de moderne Straffesystemer ere
vist med Rette overbeviste om, at det især er gennem
Arbejdet Opdragelsen skal ske.

e. Fordi vore Institutioner endnu ikke have løst
% --- p. 80 ---
\setcounter{footnote}{0}%
\opage{80}den store Opgave at hindre Forbrydelsernes, særligt Tilbagefaldenes
Tiltagen, derfor kan man ikke slutte, at der
virkeligt skulde gives absolut uforbederlige Forbrydere,
Individer, hos hvem ingen som helst etisk Stræben skulde
kunne vækkes, og som aldeles ikke skulde kunne bringes
i Forhold til det etiske Ideal. Selv om en fuldstændig
Tilbagevenden til det menneskelige Samfundsliv ikke
skulde være mulig, kunde der dog vækkes en Drift, der
førte fremad i Retning af Idealet, var det end kun med
smaa Skridt. En Forsoning vilde da være mulig, og den
\emph{absolute} Elimination ikke nødvendig. --- Der er ingen
Grund til at tro, at Udviklingen af Straffevæsenet skulde
have naat sit Højdepunkt. Reformen er kun i sin første
Begyndelse. Man har hidtil været mere opfindsom i at
udtænke Straffe, som kunde tilfredsstille Gengældelsestrangen
og Lysten til at afskrække, end i at udfinde
Midler og Veje til at paavirke Forbrydernes Karakter
for gennem denne forsinkede Opdragelse at genoprette
hvad der formedelst Samfundets Ordning var forsømt paa
et tidligere Trin. Dette kom af den udvortes Maade, paa
hvilken man tidligere betragtede Forbrydelsen. Man holdt
sig \emph{til, at} den var begaaet, og lagde mindre Vægt paa,
\emph{hvorledes} og \emph{hvorfor} den blev begaaet. Den italienske
Skole vil ved sine Arbejder, der fra saa mange Sider
belyse Forbryderkarakteren, bidrage meget til fuldstændigt
at gøre Ende paa denne urigtige Betragtningsmaade.
Denne Fortjeneste er uafhængig af Gyldigheden af de
Slutninger, Skolen selv uddrager af sine Undersøgelser.

Det staar da efter alt dette endnu som et aabent
Spørgsmaal, om der gives menneskelige Individer, som
ere aldeles uforbederlige, saaledes at den etiske Lov, selv
om den blev nok saa individualiseret, ikke skulde kunne
komme til at gøre sig gældende hos dem. For den almindelige
Teori, jeg har villet gøre gældende, er dette
% --- p. 81 ---
\setcounter{footnote}{0}%
\opage{81}Spørgsmaals Besvarelse ikke af afgørende Betydning. Men
hvis min Teori er rigtig, opfordrer den til stor Varsomhed,
inden man erklærer Individer for \emph{absolut} uimodtagelige
i etisk Henseende. Dersom hvert Individ har
sin egen individuelle Lov, er hint Spørgsmaal langt vanskeligere
at besvare, end hvis man havde Ret til at gaa ud
fra en objektiv Lov, der skulde kunne fyldestgøres af Alle.
Der udfordres en Sans for individuelle Forskelligheder og
en Taalmodighed, som giver sig Tid, indtil den haarde
Skal kan sønderbrydes. Overfor Forbryderne (»Forbrydertypen«)
har man længe baaret sig ad, lige som Europæerne
pleje at bære sig ad over for de andre Racer.
Lige som man uden videre har erklæret disse for at
mangle Evne til Civilisation\footnote[1]{De Betragtninger, med hvilke Th. \emph{Waitz} imødegaar Paastanden om Indianernes Mangel paa Evne til Civilisation, have i denne Sammenhæng Interesse. Han siger bl. a.: „Indianernes lange Stillestaaen paa et meget lavt Kulturtrin, beviser Lidet eller Intet mod deres Begavelse, saa længe de endnu ikke længere Tid igennem have været stedte i saadanne Forhold, som over Hovedet gøre et Fremskridt muligt, og det er hidtil ikke sket. Ja, endnu mere. Naar der mange Aarhundreder igennem, saa vidt vi vide, ikke har indvirket Noget paa et Folk, af hvilket man kunde love sig en omdannende Indflydelse paa deres Organisation og deres Levevis, paa deres ydre og indre Liv, saa er det sandsynligt, at der behøves langt stærkere Magter og en i langt længere Tid uafbrudt fortsat Paavirkning, for i det Hele at frembringe en Omdannelse af dem, selv en nok saa ubetydelig, end i det modsatte Tilfælde“. Die Indianer Nordamerikas. Leipzig 1865. p. 74.} og dermed dømt dem til
Undergang, saaledes har man erklæret hine for at mangle
moralsk Evne og dermed dømt dem til »Elimination«.

6. Til Slutning vil jeg lidt videre udvikle den i det
Foregaaende udtalte Tanke, at de individuelle Forskelligheders
Betydning i Etiken ikke er udtømt ved, at de
hidføre Nødvendigheden af en gennemgaaende Individu%
% --- p. 82 ---
\setcounter{footnote}{0}%
\opage{82}alisering af den etiske Lov, eller derved, at Livet faar
større Fylde og Rigdom, end hvis det skulde og kunde
tilskæres efter visse bestemte, almindelige Regler. Anerkendelsen
af dens Betydning gør tillige Slægtens Fremadskriden
forstaaelig. --- Udviklingen inden for den orgaganiske % PRINTED AS IS (orgaganiske: orga-|ganiske at the line turn (600 ppi))
Verden bliver, som Darwin har lært, kun mulig
ved, at der finder en uafbrudt Varieren Sted. Af de
»spontane« individuelle Variationer udvælges og fastholdes
de, som ere mest tjenlige for Livets fortsatte Bestaaen.
Disse saaledes fastholdte Variationer blive altsaa typiske
for Livets fortsatte Udvikling. En Analogi hertil have
vi paa det etiske Omraade. De, hos hvem de Karakteregenskaber, den etiske Lov fordrer, udvikle sig med
største Kraft og i den højeste Grad, og af hvem den
indre og ydre Modstand overvindes med Lethed eller,
selv om det sker gennem haard Kamp, paa en saadan
Maade, at de derved gennembryde Skranker, der hidtil
stansede Udviklingen, --- de staa vel i en vis Forstand
ikke etisk højere end de Mange, der maa kæmpe haardt
for blot at naa et Jevnmaal af Udvikling, og som væsentligt
holde sig paa de hidtil banede Veje, men de staa dog
som Forbilleder, der vejlede, opmuntre og trøste; de angive
ikke blot \emph{Vej og Retning}, som Mønstre til Efterfølgelse,
men vise tillige, \emph{hvilke Muligheder der
indeholdes i den menneskelige} Natur. Trods de
store individuelle Forskelligheder ere dog alle menneskelige
Væsener beslægtede, og det er en naturlig Slutning,
at hvad de videst Komne have naat, ogsaa til en vis
 Grad maa kunne blive tilgængeligt for de Andre. Som
\emph{Herbert Spencer} har sagt: »Hvad der nu karakteriserer
de undtagelsesvis Store kan ventes en Gang at
ville karakterisere Alle. Ti hvad den bedste menneskelige
Natur er i Stand til, ligger inden for den hele menneske%
% --- p. 83 ---
\setcounter{footnote}{0}%
\opage{83}
lige Naturs Rækkevidde«\footnote[1]{Data of Ethics. p. 256 f.}. Historien viser, at de Mennesker,
der staa som etiske Forbilleder for lange Tider, ikke
altid vare saadanne, hvis Indre var fuldt af stærke og
varige indre Modsætninger, eller hvis Karakter bar Spor
af store indre Kampe, men snarere harmoniske Naturer,
som levede og aandede i store Tanker, under hvilke de
andre Dele af deres Natur saa godt som uden Modstand
bøjede sig. Og det er naturligt nok, ti den, hvis Kraft
beslaglægges af den indre Kamp, vil ikke let for Andre
kunne frembyde Billedet af en mod et betydningsfuldt
Maal rettet kraftig personlig Virksomhed. I den religiøse
Forestilling er dette endog gennemført i en saadan Grad,
at man har gjort Forbilledet syndefrit og villet tage ethvert
Spor af virkelig, menneskelig Kamp og Stræben
bort. Dette er at skyde forbi Maalet; men det vidner
om Trangen til ophøjede Forbilleder, en Trang, som ikke
kan falde bort, naar Livet ikke skal blive fladt og sløvt.

Slægtens Solidaritet ophæves altsaa ikke ved den
etiske Lovs individuelle Relativitet. Slægtens Enhed bestaar
netop i og gennem Mangfoldigheden af Individer.
Det er, som jeg har søgt at vise, i Livets fælles Retning
og Tendents, at Berettigelsen til at tale om etisk Lov
begrundes. Og derved godtgøres ogsaa Slægtens Enhed
trods alle Forskelligheder.
\end{document}
